\documentclass[output=paper,colorlinks,citecolor=brown]{langscibook}
\author{Björn Wiemer\affiliation{Johannes-Gutenberg-Universität Mainz}}
\title[Two major apprehensional strategies in Slavic]
      {Two major apprehensional strategies in Slavic: A survey of their areal and grammatical distribution}
\abstract{The paper gives a first comprehensive survey of two types of apprehensional strategies which are widely attested in Slavic languages. The first type foregrounds the speaker's apprehension that some undesirable event may occur (or might have occurred), the second one builds on the first one: the speaker, assuming that an undesirable event is imminent, admonishes the addressee to undertake (or to abandon) some action in order to prevent that undesirable event. The latter type is known as preventive and most commonly expressed by perfective verbs in the negated imperative (or equivalent clause types). Since this strategy is a variety of directive speech acts, it is here called \textsc{negative-directive}. The former variety is here dubbed \textsc{negative-optative} strategy, since it is marked by clause-initial connectives most of which are used as negated optatives, but it subsumes also patterns associated with negative purpose clauses. This article shows how both strategy types complement each other conceptually, how variable the means of expression are and to which extent their components have lost compositionality, how the patterns are distributed over Slavic languages, which ones can occur in clausal complements and whether usage in independent clauses can be considered a result of insubordination.}

%move the following commands to the "local..." files of the master project when integrating this chapter
%\usepackage{tabularx}
%\usepackage{langsci-optional}
%\usepackage{langsci-gb4e}




\IfFileExists{../localcommands.tex}{
 \addbibresource{../localbibliography.bib}
 \usepackage{tabularx,multicol}
\usepackage{url}
\urlstyle{same}

\usepackage{langsci-optional}
\usepackage{langsci-lgr}
\usepackage{langsci-gb4e}

\babelprovide[import]{hebrew}

\usepackage{paralist} %%may not be needed, currently used in questionnaire


%Siena added packages for introduction; not sure they are needed
\usepackage[normalem]{ulem}
\useunder{\uline}{\ul}{}

%from Treis:
\usepackage{listings}
\lstset{basicstyle=\ttfamily,tabsize=2,breaklines=true}

%%from Francez
% % % \usepackage{cjhebrew}
%\usepackage{tipa}  Also included by D&A, need to check with LSP whether to change it to language science tipa
\usepackage{leipzig}
%%from Dabkowski&AnderBois:

\usepackage[shortlabels]{enumitem} 
\usepackage{centernot}
\usepackage{arydshln}
\usepackage{newunicodechar}
\usepackage{csquotes}
%\let\sups\relax \usepackage{tipa} commented out on 7/3/25 Martina
% \usepackage[all]{nowidow}
\usepackage{listings} \lstset{basicstyle=\ttfamily}
\usepackage{multirow}
\usepackage{rotating}
\usepackage{pdflscape}
\usepackage{xspace}
%\usepackage{tabularx}
%\usepackage{multicol}
\usepackage{supertabular}
\usepackage{hhline}
\usepackage{bbding} %this package is for the checkmarks in Table 6 in Browne et al.

%\usepackage{qtree}
%\usepackage{tree-dvips}

\usepackage{array}
\usetikzlibrary{calc,tikzmark,matrix}
\usepgflibrary{shadings}
\usepackage{pifont}
%from Fedotov:

\usepackage{langsci-textipa}
%\usepackage{multirow}


%\usepackage{xeCJK}

%%%%%%%%%%%%%%%%%%%%%%%%%%%%%%%%%%%%%%%%%%%%%%%%%%%%
%%% EXAMPLES %%%%%%%%%%%%%%%%%%%%%%%%%%%%%%%%%%%%%%%
%%%%%%%%%%%%%%%%%%%%%%%%%%%%%%%%%%%%%%%%%%%%%%%%%%%% 

% if you want the source line of examples to be in
% italics, uncomment the following line
\renewcommand{\exfont}{\itshape}
% to add additional information to the right of
% examples, uncomment the following line

%\usepackage{langsci/styles/jambox}


%%% FOOTNOTE EXAMPLE NUMBERING
% \makeatletter
% \patchcmd{\@footnotetext}{\setcounter{fnx}{0}}{\renewcommand{\thexnumi}{\roman{xnumi}}}{}{}
% \apptocmd{\@footnotetext}{
%     \@noftnotetrue
%     \renewcommand{\thexnumi}{\arabic{xnumi}}%
% }{}{}
%
% \makeatother
\usepackage[linguistics,edges]{forest}
\usepackage{hyphenat}
\usepackage{langsci-branding}

 \makeatletter
\let\thetitle\@title
\let\theauthor\@author
\makeatother

\newcommand{\togglepaper}[1][0]{
   \bibliography{../localbibliography}
   \papernote{\scriptsize\normalfont
     \theauthor.
     \titleTemp.
    To appear in:
    Martina Faller,  Marine Vuillermet \&  Eva Schultze-Berndt (eds.). Apprehensional Constructions in a cross-linguistic perspective. Berlin: Language Science Press. [preliminary page numbering]
   }
   \pagenumbering{roman}
   \setcounter{chapter}{#1}
   \addtocounter{chapter}{-1}
}

\newbool{bookcompile}
\booltrue{bookcompile}
\newcommand{\bookorchapter}[2]{\ifbool{bookcompile}{#1}{#2}}


\newfontfamily\FedotovBoldEmptySetFont{LibertinusSerif-Italic.otf}[AutoFakeBold=2.5]
\newcommand{\FedotovBoldEmpty}{{\FedotovBoldEmptySetFont\bfseries ∅}}

%From Dabkowski&AnderBois


\let\sc\scshape
\let\rm\rmfamily
%\let\it\itshape
\let\tt\ttfamily
\let\nr\normalfont

\newcommand{\wsc}[1]{\textls[80]{\scshape#1}}

\let\textai\textit
% \let\textlc\spacedlowsmallcaps
% \let\textac\spacedallcaps
\let\textlc\textsc
\let\textac\textsc


%%%%%%%%%%%%%%%%%%%%%%%%%%%%%%%%%%%%%%%%%%%%%%%%%%%%
%%% FLAG %%%%%%%%%%%%%%%%%%%%%%%%%%%%%%%%%%%%%%%%%%%
%%%%%%%%%%%%%%%%%%%%%%%%%%%%%%%%%%%%%%%%%%%%%%%%%%%%

% \newcommand{\scott}  [1]{\textcolor{violet}{#1}}
% \newcommand{\maks}   [1]{\textcolor{blue}{#1}}
% % \newcommand{\data}   [1]{\textcolor{teal}{#1}}
% \newcommand{\new}   [1]{#1}
% \newcommand{\newnew}   [1]{\textcolor{teal}{#1}}
% \newcommand{\rev}   [1]{\textcolor{magenta}{#1}}



%%%%%%%%%%%%%%%%%%%%%%%%%%%%%%%%%%%%%%%%%%%%%%%%%%%%
%%% TABLES %%%%%%%%%%%%%%%%%%%%%%%%%%%%%%%%%%%%%%%%%
%%%%%%%%%%%%%%%%%%%%%%%%%%%%%%%%%%%%%%%%%%%%%%%%%%%%

%\newcolumntype{Y}{>{\arraybackslash}p{25.3543464286pt}} % length of L

% \newcolumntype{K}{>{\arraybackslash}p{24.25pt}}
% \newcolumntype{V}{>{\arraybackslash}p{(\textwidth/9-24pt)/2}}
%
\let\ch\relax \newcommand{\ch}[2]{\multicolumn{#1}{c}{\sc#2}}
\let\sx\relax \newcommand{\sx}[2]{\multicolumn{#1}{l}{\emph{-#2}}} % suffix
\let\su\relax \newcommand{\su}{\sx{1}} % suffix unicolumnal
\let\sd\relax \newcommand{\sd}{\sx{2}} % suffix dicolumnal
\let\cx\relax \newcommand{\cx}[2]{\multicolumn{#1}{l}{\emph{=#2}}} % clitic
\let\cu\relax \newcommand{\cu}{\cx{1}} % clitic unicolumnal
\let\cd\relax \newcommand{\cd}{\cx{2}} % suffix dicolumnal
\let\gx\relax \newcommand{\gx}[2]{\multicolumn{#1}{l}{\sc #2}} % gloss
\let\gd\relax \newcommand{\gd}{\gx{2}} % gloss dicolumnal
\let\gr\relax \newcommand{\gr}[1]{\multicolumn{2}{c}{\multirow{2}{*}{\textcolor{gray}{\sc#1}}}}
\let\db\relax \newcommand{\db}[1]{\multicolumn{2}{c}{\multirow{2}{*}{\sc#1}}}



%%%%%%%%%%%%%%%%%%%%%%%%%%%%%%%%%%%%%%%%%%%%%%%%%%%%
%%% EXAMPLES %%%%%%%%%%%%%%%%%%%%%%%%%%%%%%%%%%%%%%%
%%%%%%%%%%%%%%%%%%%%%%%%%%%%%%%%%%%%%%%%%%%%%%%%%%%%

%%% AFFIX BOUNDARY %%%%%%
% \DeclareRobustCommand{\longmn}{\textminus}
% \DeclareRobustCommand{\mn}{-}

%%% CLITIC BOUNDARY %%%%%%
% \DeclareRobustCommand{\longeq}{=}
% \makeatletter\DeclareRobustCommand{\eq}{%
%     \settowidth{\@tempdima}{-}% width of a hyphen
%     \resizebox{\@tempdima}{\height}{=}%
% }\makeatother

%%% FUZZY CLITIC %%%%%%
% \DeclareRobustCommand{\longfc}{%
%     $\overset{\text{\smash{}}}{\text{$\approx$}}$%
% }
% \makeatletter\DeclareRobustCommand{\fc}{%
%     \settowidth{\@tempdima}{-}% width of a hyphen
%     \resizebox{\@tempdima}{\height}{%
%         $\overset{\text{\smash{}}}{\text{$\approx$}}$%
%     }%
% }\makeatother

%%% REDUPLICATION %%%%%%%
% \DeclareRobustCommand{\longtl}{$\sim$}
% \makeatletter\DeclareRobustCommand{\tl}{%
%     \settowidth{\@tempdima}{-}% width of a hyphen
%     \resizebox{\@tempdima}{\height}{$\sim$}%
% }\makeatother

% \newcommand{\mn}          {\textminus}
\newcommand{\src}    [1]  {\jambox[0pt]{\rm(#1)}}
% \newcommand{\ai}     [2]  {\emph{#1} `\textsc{#2}#3'\xspace}
\newcommand{\ai}     [2]  {\emph{#1} `\textsc{#2}'\xspace}
\newcommand{\sane}   [1][]{\ai{=sa'ne}{appr}{#1}}
\newcommand{\srcn}   [1][]{\ai{kûintsû}{srcn}{#1}}
% \newcommand{\mbekane}[1][]{\ai{\eq mbe kañe}{neg.adv aux.inf}{#1}}

% \let\sc\relax \newcommand{\sc}{\normalfont\scshape}
% \let\rm\relax \newcommand{\rm}{\normalfont\rmfamily}
% \let\it\relax \newcommand{\it}{\normalfont\itfamily}

\newcommand{\lb}{{\rm[}}
\newcommand{\rb}{{\rm]}}

\newcommand{\bad}{\makebox[0pt][r]{*\,}}
\newcommand{\infel}{\makebox[0pt][r]{\#\,}}


% \newcommand{\lsqz}[1]{#1}
% \newcommand{\lsqz}[1]{\makebox[0pt][l]{#1}}

% \newcommand{\lsqu}[1]{\makebox[0pt][l]{#1}}
% \newcommand{\rsqu}[1]{\makebox[0pt][r]{#1}}
% \newcommand{\astr}{\makebox[0pt][r]{{\nr*}}}
% \newcommand{\hash}{\makebox[0pt][r]{{\nr\#}}}
% \newcommand{\qstn}{\makebox[0pt][r]{{\nr?}}}
% \newcommand{\bad}{\makebox[0pt][r]{*}}
% \newcommand{\unhappy}{\makebox[0pt][r]{\#}}


%%%%%%%%%%%%%%%%%%%%%%%%%%%%%%%%%%%%%%%%%%%%%%%%%%%%
%%% UNICODE CHARACTERS  %%%%%%%%%%%%%%%%%%%%%%%%%%%%
%%%%%%%%%%%%%%%%%%%%%%%%%%%%%%%%%%%%%%%%%%%%%%%%%%%%

\newunicodechar{⟨}{⟨}
\newunicodechar{⟩}{⟩}

\newcommand{\infix}{⟨}
\newcommand{\outfix}{⟩}



%%%%%%%%%%%%%%%%%%%%%%%%%%%%%%%%%%%%%%%%%%%%%%%%%%%%
%%% IPA %%%%%%%%%%%%%%%%%%%%%%%%%%%%%%%%%%%%%%%%%%%%
%%%%%%%%%%%%%%%%%%%%%%%%%%%%%%%%%%%%%%%%%%%%%%%%%%%%

% \newcommand{\ipa}{\textipa}

\newcommand{\sub}{\textsubscript}
% \let\super\relax \newcommand{\super}{\textsuperscript}

\newcommand{\ts}{\texttslig}
% \let\dz\relax \newcommand{\dz}{\textdzlig}
\newcommand{\tS}{\textteshlig}
\newcommand{\dZ}{\textdyoghlig}
\newcommand{\cj}{\textctc}
\newcommand{\sj}{\textctc}
\newcommand{\zj}{\textctz}
\newcommand{\tcj}{\texttctclig}
\newcommand{\tsj}{\texttctclig}
\newcommand{\dzj}{\textdctzlig}
\newcommand{\gj}{\textbardotlessj}
% \let\nj\relax \newcommand{\nj}{\textltailn}
% \newcommand{\W}{\textturnmrleg}
\newcommand{\wl}{\textturnmrleg}
% \newcommand{\1}{\ipabar{\~\i}{.5ex}{1.1}{}{}}
\newcommand{\dotlessibar}{\ipabar{\i}{.5ex}{1.1}{}{}}
\newcommand{\darkl}{\textltilde}
% \newcommand{\n}{\textsubarch}
% \newcommand{\x}{\textovercross}
\newcommand{\lowered}{\textlowering}
\newcommand{\raised}{\textraising}
\newcommand{\retracted}{\textsubbar}
\newcommand{\unreleased}{\textcorner}
% \newcommand{\syllabic}{\s}
\newcommand{\implosivet}{\texthtt}
\newcommand{\implosived}{\texthtd}
\newcommand{\implosivep}{\texthtp}
\newcommand{\implosiveb}{\texthtb}
\newcommand{\implosivek}{\texthtk}
\newcommand{\implosiveg}{\texthtg}



%%%%%%%%%%%%%%%%%%%%%%%%%%%%%%%%%%%%%%%%%%%%%%%%%%%%
%%% OTHER %%%%%%%%%%%%%%%%%%%%%%%%%%%%%%
%%%%%%%%%%%%%%%%%%%%%%%%%%%%%%%%%%%%%%%%%%%%%%%%%%%%

\newcommand{\ie}{i.\,e.\xspace}
\newcommand{\Ie}{I.\,e.\xspace}
\newcommand{\eg}{e.\,g.\xspace}
\newcommand{\Eg}{E.\,g.\xspace} 




\newcommand{\francoissource}[1]{\hfill\textnormal{{\footnotesize #1} }}
\newcommand{\E}{Ɛ\kern-1.5pt\raisebox{1pt}{̏}\kern1.5pt}
\newcommand{\ù}{{ṵ}\kern-2pt{̀}\kern2pt}
\newcommand{\ȕ}{{ṵ}\kern-2pt{̏}\kern2pt}

\newcommand{\OO}{ɔ\kern-1pt\raisebox{-1pt}{̰}\kern1pt}
\newcommand{\Ǒ}{{{ɔ̌\kern-1pt\raisebox{-.5pt}{̰}\kern1pt}}}


\newcommand{\à}{{a̰}\kern-2pt{̀}\kern2pt}
\newcommand{\ilit}[1]{#1\il{#1}}

 %% hyphenation points for line breaks
%% Normally, automatic hyphenation in LaTeX is very good
%% If a word is mis-hyphenated, add it to this file
%%
%% add information to TeX file before \begin{document} with:
%% %% hyphenation points for line breaks
%% Normally, automatic hyphenation in LaTeX is very good
%% If a word is mis-hyphenated, add it to this file
%%
%% add information to TeX file before \begin{document} with:
%% %% hyphenation points for line breaks
%% Normally, automatic hyphenation in LaTeX is very good
%% If a word is mis-hyphenated, add it to this file
%%
%% add information to TeX file before \begin{document} with:
%% \include{localhyphenation}

\hyphenation{
par-a-digm
com-ple-men-ta-tion
anaph-o-ra
Dor-drecht
Cu-shi-tic
be-ne-dic-tive
pa-ra-digms
Ethi-o-pi-an
hoch-land-ost-ku-schi-ti-schen 
Duu-raa-me
Pa-ken-dorf
Lich-ten-berk
affri-ca-te
affri-ca-tes
com-ple-ments
Eng-lish
}




\hyphenation{
par-a-digm
com-ple-men-ta-tion
anaph-o-ra
Dor-drecht
Cu-shi-tic
be-ne-dic-tive
pa-ra-digms
Ethi-o-pi-an
hoch-land-ost-ku-schi-ti-schen 
Duu-raa-me
Pa-ken-dorf
Lich-ten-berk
affri-ca-te
affri-ca-tes
com-ple-ments
Eng-lish
}




\hyphenation{
par-a-digm
com-ple-men-ta-tion
anaph-o-ra
Dor-drecht
Cu-shi-tic
be-ne-dic-tive
pa-ra-digms
Ethi-o-pi-an
hoch-land-ost-ku-schi-ti-schen 
Duu-raa-me
Pa-ken-dorf
Lich-ten-berk
affri-ca-te
affri-ca-tes
com-ple-ments
Eng-lish
}



 \boolfalse{bookcompile}
 \togglepaper[23]%%chapternumber
}{}






%\newcommand{\tikzmark}[1]{\tikz[remember picture] \node (#1) {};}

\begin{document}
\maketitle 


\section{Introduction}
\label{sec:intro}

Slavic languages lack dedicated morphosyntactic markers of apprehensional meanings, but they have at their disposal constructions which can be regarded as apprehensional strategies, i.e.\ as patterns that are regularly associated with apprehensional semantics.\footnote{I will employ \textit{apprehensive} only with reference to dedicated markers and use \textit{apprehensional} only for the semantics. For a brief terminological overview cf.\ \citet[59]{MitkovskaEtAl2017}.} Apart from them, there are various lexical means, usually referred to as particles.\footnote{Cf.\ \citet[46--47]{Dobrushina2006wiemer}, \citet{Arkadiev2007}, \citet[448]{Plungian2011}, \citet{Mikaelian2021} for Russian and \citet{Ivanova2014} for Russian and Bulgarian.} Many of them are employed as clause-initial connectives, with different degrees of syntactic integration on a coordination-subordination cline, but there is only one strategy, the negative optative strategy, which can clearly be considered as part of clausal complementation after suitable predicates denoting apprehension (see \sectref{subsec:insubordination} below).

Here I will focus on two strategies which are widespread across the Slavic landscape, which can occur in isolated utterances and which complement each other conceptually. One of them is derived from negative optative or negative purpose clauses (via extension into epistemic modality) (see \ref{ex:wiemer1}--\ref{ex:wiemer2}), the other from imperatives or equivalent constructions coding directive speech acts (see \ref{ex:wiemer3}--\ref{ex:wiemer4}). \\

\noindent \textbf{Negative-optative strategy}

\ea\label{ex:wiemer1} \textnormal{Russian} \\
\gll \textbf{Kak} \textbf{by}\textnormal{/}\textbf{Liš'} \textbf{by} ego ne oštrafova-l-i. \\
how \textsc{irr}/only \textsc{irr} \textsc{3sg.m.acc} \textsc{neg} fine\textsc{[pfv]-pst-pl} \\
\glt `[I'm afraid] they might fine him.'
\z 


\ea\label{ex:wiemer2} \textnormal{Bulgarian (\citealt[147]{Ivanova2014})}\\
\gll T-oj približi s nebrežn-o izraženi-e na lic-e-to -- \textbf{da} \textbf{ne} \textbf{bi} tija, malki-te, da si pomislj-at nešto. \\
3-\textsc{nom.sg.m} approach\textsc{[pfv].aor.3.sg} with indifferent-\textsc{sg.n} expression\textsc{[n]-sg} on face-\textsc{sg.n-def.sg.n} {} \textsc{irr} \textsc{neg} \textsc{sbjv} \textsc{3pl} small-\textsc{def.pl} \textsc{irr} \textsc{refl.dat} think\textsc{[pfv]-prs.3.pl} something \\
\glt `He came closer with an indifferent expression on his face -- lest the little ones think something (uncomfortable for him).'
\z

\noindent \textbf{Negative-directive strategy}
\ea\label{ex:wiemer3}
\textnormal{Russian} \\
\gll (Smotr-i,) ne poskol'zn-i-s'! \\
look\textsc{[ipfv]-imp.sg} \textsc{neg} slip\textsc{[pfv]-imp.sg-refl} \\
\glt `(Be careful,) don't slip!'
\z

\ea \textnormal{Bulgarian} \\
\label{ex:wiemer4}
\gll (Pazi se) da ne padn-eš! \\
look\textsc{[ipfv]:imp.sg} \textsc{refl} \textsc{irr} \textsc{neg} fall\textsc{[pfv]-prs.2sg} \\
\glt `(Look) that you don't fall!'
\z

 The first pattern primarily marks just apprehension, i.e.\ the speaker evaluates a situation both as likely to occur (or have occurred) and as undesirable.\footnote{This is probably the most widespread understanding of the terms \textit{apprehensive/apprehensional} in Slavic Linguistics. Many (e.g., \citealt{Ivanova2014}, \citealt{Mitkovska2015}, \citealt{Bužarovska2017}, \citealt{MitkovskaEtAl2017}) restrict them to this type. Typically, these researchers focus on Balkan Slavic, while researchers dealing with the negative-directive strategy (see next footnote) have focused on Russian (and other North Slavic languages).} From a conceptual point of view, the second pattern builds on the first one, but encourages the addressee to do (or to abandon) something in order to prevent this undesirable situation. Correspondingly, this pattern has been labelled \textit{preventive}, and it has been treated in connection with directive speech acts.\footnote{Cf.\ \citet{Laskowski1998}, \citet{Xrakovskij2001}, \citet{Nilsson2012,Nilsson2013},  \citet{Kissine2014}, to name but a few.} For the sake of brevity I will call the latter pattern \textsc{NegDir} and the former \textsc{NegOpt}. The label \textsc{NegOpt} comprises constructions related either to optative or to purpose clauses. This joint treatment of negative optative and negative purpose sources seems justified because in their contemporary stages they are more similar to each other in structure and distribution than either of them to \textsc{NegDir}. This also holds true for their provenance, even though in comparison to \textsc{NegDir} (\sectref{sec:negdir}) the relevant cases of \textsc{NegOpt} and their diachronic sources are more diversified (\sectref{sec:negopt} and \sectref{sec:diachronic}). If not stated otherwise, claims concerning \textsc{NegOpt} will pertain to structures irrespective of their provenance from properly negative-optative or from negative-purpose clauses. Both \textsc{NegDir} and \textsc{NegOpt} constructions occur as main clauses; \textsc{NegDir} does so exclusively, while \textsc{NegOpt} can occur embedded as a clausal complement. Either of them can also be uttered in isolation, i.e.\ out of the blue in dialogue, whereas most of the less salient apprehensional strategies cannot.


As for the latter, it is worth emphasizing that at least in some Slavic languages we find strategies which are based on metalinguistic justification. These strategies cannot be uttered in isolation, as they depend on preceding verbal discourse. A prominent case in point is the Polish causal conjunction \textit{bo} `because' (cf.\ \citealt{Kawka1988}), illustrated in (\ref{ex:wiemer5}).

\largerpage[2]
\ea \label{ex:wiemer5}
\langinfo{Polish}{}{\citealt[135]{Ożóg1990}}\\
\gll Nie rusza-j t-ego \textbf{bo} się oparz-ysz. \\
\textsc{neg} touch[\textsc{ipfv}]-\textsc{imp.sg} \textsc{dem-gen.sg} because \textsc{refl} burn\textsc{[pfv]-fut.2sg} \\
\glt `Don't touch this, \textbf{or} [lit.\ `since'] you'll burn yourself.' ($\cong$ `$\ldots$ lest you burn yourself.') \\
\z

\noindent The causal connective introduces the second conjunct in a clause chain; the preceding conjunct usually expresses a directive speech act. As a rule, the predicate of this conjunct with the causal connective is a perfective verb in the present tense stem (by default interpreted as future). In any case, the conjunct following the directive speech act is interpreted as a warning (or admonition) and verbalizes the consequences that are deemed to follow if the addressee does not comply with this speech act. Such consequences can include concerns for the interlocutor's benefit, in which case the expression acquires an apprehensional function.


Here causal and equivalent connectives\footnote{Functional equivalents are metadiscursive uses of the Bulgarian conjunction \textit{če} (today the standard \textsc{that}-complementizer, but deriving from a causal function) and the Russian adversative connective \textit{a to}. All these connectives have other functions outside the domain of apprehension and their apprehensional interpretation is highly dependent on the discourse context.} will be left out of consideration. The motivation to concentrate on the aforementioned apprehensional constructions, \textsc{NegOpt} and \textsc{NegDir}, is as follows. First, to my knowledge, there has not been so far any comprehensive account of apprehensional strategies in Slavic at all. To start with (especially in view of space restrictions), it seems justified to choose those constructions which are most salient in terms of areal spread, conventionalization and syntactic independence, and for which there already exists enough reliable research. Second, these constructions show an interesting areal distribution which can be described in a more fine-grained manner if we take into account their dependence on core distinctions in the grammar of Slavic languages (TAM distinctions). Moreover, subareas of Slavic differ considerably in the degree with which they distinguish these two strategies in structural and functional terms. Third, only for the two strategies chosen can we supply some limited diachronic background, whereas for the others hardly any such research has been carried out (at least I am not aware of any), whether corpus-based or not. One of the reasons certainly lies in the fact that apprehensional strategies are often difficult to discern, since they are tightly bound to discourse pragmatics, and their degree of conventionalization is low. Jointly with this, and fourth, I want to avoid discussions concerning the coordination-subordination gradient, which would arise for many of the clause-initial particles mentioned in the literature. In my conviction, discussions about where such structures are to be located on a coordination-subordination cline do not yield much for an understanding of apprehensional strategies. The \textsc{NegOpt} strategy is probably the only one in Slavic languages which really can occur as a syntactically embedded structure. For this reason I will pay attention to the issue of insubordination in connection with \textsc{NegOpt} (see \sectref{subsec:insubordination}).


The article has the following structure. I will first discuss conceptual matters important to the understanding of how the two main strategies are related on the semantics-pragmatics interface. For this purpose I will introduce templates (\sectref{sec:conprem}). Subsequently, I will analyze the two patterns one by one, first \textsc{NegOpt} (\sectref{sec:negopt}), then \textsc{NegDir} (\sectref{sec:negdir}). Each of these sections will account for the areal spread and differentiation in the contemporary languages. This account will be somewhat uneven, because for many Slavic languages sufficiently reliable data are lacking or not easy to access. In particular, I will not have anything to say here on Ukrainian and Belarusian, nor on Sorbian. In \sectref{sec:diachronic} I will discuss some diachronic background, in particular the issue of insubordination, which is relevant only for \textsc{NegOpt} (\sectref{subsec:insubordination}). The last section gives a summary of findings and provides an outlook (\sectref{sec:summary}).


I will be referring to subareas of Slavic following commonplace divisions into East, West and South Slavic, as well as Balkan Slavic. North Slavic unites West and East Slavic. Balkan Slavic comprises Bulgarian, Macedonian and Torlak and Timok dialects in southeastern Serbia. I will gloss clause-initial connectives according to their etymological sources; however the reader should have in mind that as apprehensional markers (or in similar functions) they have lost transparency (see \sectref{sec:negopt}). \textsc{con} is used to gloss connectives if no further specification is warranted. Examples given in the running text will be glossed to a minimum, aspect will then be marked by upper case indexes (\textsc{\textsuperscript{pfv}}, \textsc{\textsuperscript{ipfv}}).


\section{Conceptual premises} \label{sec:conprem}

Apprehensional strategies can be classified in various ways. From a conceptual point of view, some researchers (e.g., \citealt{Dobrushina2006wiemer}, \citealt{Ivanova2014}) regard the \textsc{NegOpt} pattern as basic and the \textsc{NegDir} pattern as secondary (see Footnote 3). Other researchers take the opposite view, since they approach apprehensional semantics from imperatives and directive speech acts (e.g., \citealt{Birjulin1992,Birjulin1994}; \citealt{Aikhenvald2010}; \citealt{Gusev2013}). After all, apprehensives and apprehensional strategies are conceptually very complex. When descriptions by various scholars are compared,\footnote{Cf.\ \citet[293--296]{lichtenberk1995}; \citeauthor{Dobrushina2006wiemer} (\citeyear[30, 49]{Dobrushina2006wiemer}, \citeyear{dobrushina2021}), \citet[448--449]{Plungian2011}, \citet[258--259]{angeloschultzeberndt2016}.} three components belonging to different domains become salient. See \tabref{tab:componentsslavic}, designed to apply to the analysis of Slavic apprehensional strategies, which I focus on here. Thus, it is not meant as a general conceptual schema of apprehensive markers and/or apprehensional strategies; not all scholars distinguish, or acknowledge, all three components.

\begin{table}

\caption{Components of apprehensional strategies (at least in Slavic languages)}
\label{tab:componentsslavic}
\begin{tabularx}{\textwidth}{lQ}
\lsptoprule
1. Epistemic assessment  & The speaker lends some knowledge- or belief- based support that a certain eventuality E has obtained, is  obtaining or can be predicted:  `I~consider E likely  to occur/ to have occurred/ to be the case.' \\ \midrule
2. Undesirability  & The speaker considers E, or its consequences,  undesirable and doesn't want them to occur;  this presupposes an element of volition:  `I do not want E, or its consequences, to occur.'                    \\ \midrule
3. Uncontrollability & Neither a discourse participant (in particular the  addressee) nor a third party can influence the  course of events.                                                                                        \\
\lspbottomrule
\end{tabularx}
\end{table}

Components 1 and 2 are separated from each other in \citet[114--115]{Birjulin1992,Birjulin1994} and \citet[98]{Nilsson2013}; they are also discussed by \citet[83--88]{Holvoet2010} and \citet[62]{Gusev2013}. All these authors concentrate on preventives, but cf.\ also \citet{MitkovskaEtAl2017}. Preventives conjoin two elements: (i) an epistemic (`P may/is bound to occur') and (ii) a directive one (`do something so that P doesn't occur'). This is why Gusev characterizes them by features typical for both indicative and imperative forms.

Component 2 appears to be uncontroversial; in fact, it spells out the core of apprehensional marking. It can however be implied in a warning and, in this sense, be somewhat downgraded in communicative salience. 

\largerpage
The epistemic component (Component 1) has been questioned by some researchers discussing co-occurrence and replacement conditions of apprehensive modals vis-à-vis other modals. In particular, one may question whether apprehensive markers, especially when they are exclusively future-oriented, meet a restrictive definition of epistemic modality \citep{chapters/vuillermet, chapters/introduction}. Regardless of how apprehensives relate to modals, the requirement of unrestricted occurrence (in all tenses) applies to the \textsc{NegOpt} strategy in a conceptual perspective, and as for the \textsc{NegOpt} strategy in South Slavic there are also no grammatical restrictions concerning tense (see \sectref{sec:negopt}). More importantly, from a conceptual point of view, one cannot decompose the notion of apprehension without recurrence to epistemic attitudes (similarly to the concept `hope'), because one can only feel fear of something which one considers to be possible. Typically, the epistemic assessment relates to a situation posterior to a reference interval, but this is not necessarily the case.\footnote{Cf.\ \citeauthor{dobrushina2021}  (\citeyear{dobrushina2021}: \S2.1) on complements of `fear' verbs.}

Component 3 is likewise controversial, at least as a general feature of apprehensional patterns, but for a different reason. It has been downplayed in many discussions (e.g., \citealt[297, 312--314]{lichtenberk1995}, \citealt[95, 98--100]{Nilsson2013}, \citealt[59f., but see p. 65]{MitkovskaEtAl2017}), while other linguists have paid quite some attention (e.g., \citealt[95--96, 103--105]{Holvoet1989}, \citealt[213--215]{Xrakovskij1990}, \citealt[35--36]{Luczkow1997}). Disputes can be solved if we ask what is considered to be the object of control: the observed situation or an action to be undertaken. Moreover, it is not the speaker, but the addressee who is considered to (be able to) control that action. This becomes evident with preventives, which are realized by the \textsc{NegDir} strategy. Thus, \citeauthor{Birjulin1992} (\citeyear[5, 7]{Birjulin1992}; \citeyear{Birjulin1994}) claims that preventives imply uncontrollability as their most pervasive feature, but he then insists that this construction prompts the addressee to take measures to avert consequences (cf.\ also \citealt[225]{Aikhenvald2010}). Control(lability) is thus presupposed, although it often relates not to the event denoted by the verb, but to another, unnamed action (see \sectref{sec:negdir}). As for \textsc{NegOpt-}strategies, controllability may or may not hold (see \sectref{sec:negopt}).\footnote{Closely related `in case'-readings, in turn, are rather indifferent to both controllability \citep[58]{Dobrushina2006wiemer} and undesirability. In an `in-case' reading the speaker's assumption about a possible event need not be undesirable (e.g.\ \textit{Take a bag, in case they have food left}), and often the event is neither under the speaker's nor the addressee's control (\textit{Take an umbrella, in case it starts raining}); cf.\ \citet[298--302, 311]{lichtenberk1995}.} This settled, \citet{Gusev2013} divides the meaning of preventives into two components as shown in \tabref{tab:twocomponents}.\footnote{More precisely, an alternative should be added to subcomponent (b): `$\ldots$ or to not undertake some intermediate action that would/might lead to P' (cf.\ \citealt{Birjulin1992,Birjulin1994}, \citeauthor{BirjulinXrakovskij1992} \citeyear[38]{BirjulinXrakovskij1992}, \citeyear[37--38]{BirjulinXrakovskij2001}, \citealt[30]{Dobrushina2006wiemer}).}

\begin{table}
\caption{Two components of the meaning of preventives (\citealt[61--62]{Gusev2013}{;} my translation, BW)}
\label{tab:twocomponents}
\begin{tabular}{ll}
\lsptoprule
(a) & The speaker states that some undesirable event P may occur.\\ 
\midrule
(b) & By stating this, the speaker tries to cause the addressee or some third \\
    & person to undertake some intermediate action (P\textsubscript{aux}) which would \\
    & prevent P from occurring (i.e.\ which would cause non-P). \\
 \lspbottomrule
\end{tabular}
\end{table}

\largerpage
\citet[30, 51, 58]{Dobrushina2006wiemer} bases her typology on the fore-/backgrounded status of these two components, and this yields what I have dubbed the \textsc{NegOpt} and the \textsc{NegDir} pattern. \citet[61]{Gusev2013} makes this relation more explicit: the two constructions differ depending on which of these two components is focused. If the focus lies on component (a), i.e.\ an apprehension (= epistemic+undesirable), component (b), representing a directive illocution, remains as an implicature; under favorable circumstances it can, however, be interpreted as the proper illocutionary point. This is shown in \tabref{tab:template1} (here and in the following, grey shading indicates asserted parts of meaning or preferred interpretations, respectively). The foregrounded part can also be circumscribed as `I wish that something likely doesn't/won't happen'; for this reason, this strategy is termed the Negative Optative (\textsc{NegOpt}) strategy here.

\begin{table}

\caption{Apprehensional Template 1: Negative-optative strategy}
\label{tab:template1}
\begin{tabularx}{\textwidth}{llQ}
\lsptoprule
\textbf{Component} & \textbf{Russian} & \textbf{English} \\ \midrule
Background  & \textit{Segodnja skol'zko} & `It's slippery today.' \\
\midrule
\textbf{FOREGROUND (=a)}  & \cellcolor{lightgray}\textit{Ty mož-eš' upas-t'/upad-eš'} & `You might/will fall {[}but I don't want \\
apprehension &   &  this to happen.{]}' \\
 &  & = `If only you don't fall/\newline
 I don't wish you to fall!' \\
\midrule
\textbf{Implicature (=b)} & \textit{Xodi ostorožnee!} &  ‘Walk more carefully!’ \\
pre-emptive action; & & \\ 
can become implicit & & \\
illocutionary point \\ 
\lspbottomrule
\end{tabularx}
\end{table}

\largerpage
 If the directive component (b) is foregrounded, component (a) is presupposed. However, if the directive predicate denotes an event that is out of the addressee's control, it does not call the intermediate, pre-emptive, action (P\textsubscript{aux}) by its name; in this case P\textsubscript{aux} remains vague and the illocutionary point implicit. If, on the other hand, the directive predicate denotes an event that can be controlled by the addressee, then the illocutionary point is named directly~– compare the directive in \tabref{tab:templatetwo} (component (b)) with \textit{Oden'sja teplo!} `Put on warm clothes!' In order to tease apart these components, we have to extend Gusev's schema (\tabref{tab:twocomponents}) as in \tabref{tab:templatetwo}.\footnote{Cf.\ \citet[19]{Laskowski1998} for a similar template. \citet[15--16]{Birjulin1992} did not fully consider the possibility that the directive predicate may denote a controllable action directly, but in other respects the reasoning presented here is congruent with his view.}

\begin{table}

\caption{Apprehensional Template 2: Negative-directive strategy}
\label{tab:templatetwo}
\begin{tabularx}{\textwidth}{QQQ}
\lsptoprule
\textbf{Component}                                                                                      & \textbf{Russian}                                  & \textbf{English}                                                          \\
\midrule
Assumption                                                                                              & \textit{Na ulice xolodno}                         & `It's cold outside.'                                                      \\
\midrule
Presupposition (=a)\\Apprehension                                                                       & \textit{Prostudiš'sja. / Ty možeš' prostudit'sja} & `You'll/You might catch a cold.' \\
\midrule
FOREGROUND\newline illocutionary point (=b)\newline  directive {[}uttered{]}  & \cellcolor{lightgray}\textit{Ne prostudis'!} & `Don't catch a cold!'                                                     \\
 \midrule
\multicolumn{3}{p{11.5cm}}{If the verb in the directive denotes an uncontrollable event, illocutionary point shifts to another implied controllable action (= Gusev's P\textsubscript{aux}): } \\
& {\cellcolor{lightgray}`do something about it} \\
& {\cellcolor{lightgray}(when leaving the house)', e.g.} \\
{[}not uttered{]}     & \cellcolor{lightgray}\textit{Oden'sja teplo!} & `Put on warm clothes!'                                                    \\
\lspbottomrule
\end{tabularx}%
\end{table}


The often implicit character of the illocutionary purpose and the more complex conceptual structure of the \textsc{NegDir} pattern causes a certain asymmetry in comparison with the \textsc{NegOpt} pattern, but nonetheless both strategies complement each other in terms of their illocutionary purpose. Crucially, components (a) and (b) can be captured in this fore-backgrounding relation.

Notably, negation plays a key role in this relation. Although the “positive” counterparts of both apprehensional strategies, optatives and unnegated directives, are volition-based, they differ with respect to control; jointly with this, they differ in their association with interlocutors' roles: directives imply the addressee's control, whereas optatives rather exclude control and are often speaker-oriented (although the speaker may express concern for the addressee or some other person; see \sectref{subsec:negoptcomponents}, \sectref{subsec:negimp}). Control is maintained for negated directives and remains weak (or absent) for apprehensional expressions, but negation brings both into connection because it marks what should not happen; the difference is that in \textsc{NegOpt} strategies negation relates to the volitional evaluation (emotion), while in \textsc{NegDir} strategies it relates to what should be prevented (action).


Moreover, one gets the general impression that constructions with a negative-purpose source (see (\ref{ex:wiemer6})) seem to acquire a preventive (\textsc{NegDir}) meaning more easily than constructions with a negative-optative source (as, e.g., with Russian \textit{kak by $\ldots$  ne}; see (\ref{ex:wiemer7})). However, this remains to be investigated. Preventives are often accompanied by adverbials emphasizing that the respective event might occur inadvertently, such as \textit{czasem} `accidentally, occasionally' in (\ref{ex:wiemer6}).

\ea
\label{ex:wiemer6} 
\langinfo{Polish}{}{\citealt[182]{Danielewiczowa2014}} \\
\gll \textbf{Aby} \textbf{czasem} \textbf{nie} zdepta-l-i moich tulipan-ów! \\
\textsc{con.irr} occasionally \textsc{neg} trample.down\textsc{[pfv]-pst-pl} my tulip-\textsc{gen.pl} \\
\glt `If only they don't accidentally trample down my tulips!' \\
\z

 Furthermore, apprehensives and apprehensional strategies are inherently time-located. This is most obvious from their by far most predominant use, which relates to events that can be foreseen in the immediate future, or are even on the verge of happening (see Component 1 in \tabref{tab:componentsslavic}). Apprehensional constructions sound weird with repeated (or habitual) events. They share these two features with avertive grams, which mark that an (usually unpleasant) event was first considered imminent but nonetheless did not occur, e.g.\ \textit{I was about to fall (but didn't)} (\citealt{Kuteva1998}; \citealt{kutevaetal2019}; \citealt{Romaine2005}; however, \citealt[130]{Schmidtke-Bode2009} et passim identifies \textit{avertive} with negative purpose). Since, contrary to apprehensional constructions, avertive grams assert the non-occurrence of the unpleasant event, they are usually bound to past tense and not to perspectives posterior to speech time. Nonetheless, apprehensives and apprehensional strategies are not automatically future-oriented (\citealt{lichtenberk1995}: 296, 310, 312; \citealt{Bužarovska2017}).\footnote{\citet[17]{Plungian2004} argues that an inherent future-orientation of apprehensives explains their evaluative feature (included in Component 2 of \tabref{tab:componentsslavic}). However, there is no conceptual reason why situations that already occurred (or are occurring) should not be the object of negative evaluation, too.}

Finally, the internal temporal structure of the situation denoted by the apprehensional clause need not be bounded (or telic). Compare Russian \textit{Ja bojus', kak by on ne \textbf{spal}}\textsc{\textsuperscript{ipfv}} \textit{uže} `I'm afraid that he \textbf{is} already \textbf{asleep}' or the Macedonian example (\ref{ex:wiemer24}). In practice, this observation is of relevance only for \textsc{NegOpt}, not for \textsc{NegDir} strategies, since it hardly makes sense to ask the addressee to take care lest they accidentally perform ``a bit'' of some unbounded activity (see \sectref{sec:negdir}).


With these premises in mind, let us now deal more closely with the two strategies \textsc{NegOpt} (\sectref{sec:negopt}) and \textsc{NegDir} (\sectref{sec:negdir}).



\section{Negative-optative strategy (\textsc{NegOpt})} \label{sec:negopt}


Negative-optative strategies are related to optative or purpose clauses with different degrees of semantic transparency. I first explain the ingredients of this construction type (\sectref{subsec:negoptcomponents}), before I comment on ways of its embedding (\sectref{subsec:embedding}), on temporal relations between the content of the apprehensional clause and the time of its evaluation by the judging subject (\sectref{subsec:tempref}), and on issues concerning structural and semantic compositionality (\sectref{subsec:lossofcomp}).



\subsection{Components of the construction}
\label{subsec:negoptcomponents}

In terms of the meaning components separated in \tabref{tab:twocomponents} and in Tables \ref{tab:template1} and \ref{tab:templatetwo}, \textsc{NegOpt} strategies primarily communicate the speaker's fear that an undesirable situation might obtain (= component (a)). The prevention of this situation \mbox{(= component (b))} is not a proper part of the meaning of this strategy; it can be inferred as a sort of perlocutionary implicature \citep[29]{Dobrushina2006wiemer}, provided the relevant participant (speaker or another person) has control over some part of the situation. Compare examples (\ref{ex:wiemer1}) (`(I'm afraid) they might fine him') and (\ref{ex:wiemer7}):


\ea \textnormal{Russian}
\label{ex:wiemer7} 
\ea\label{ex:wiemer7a}
\gll Kak by\textnormal{/}Liš' by (nam) ne opozda-t'. \\
how \textsc{irr}/only \textsc{irr} \textsc{1pl.dat} \textsc{neg} come\textunderscore late\textsc{[pfv]-inf} \\
\ex\label{ex:wiemer7b}
\gll Kak by\textnormal{/}Liš' by my\textnormal{/}oni ne opozda-l-i. \\
how \textsc{irr}/only \textsc{irr} \textsc{1pl.nom}/\textsc{3pl.nom} \textsc{neg} come\textunderscore late\textsc{[pfv]-pst-pl} \\
\glt `If only we/they don't come late.'
\z
\z

In North Slavic, the construction consists of a negated verbal predicate and a clause-initial particle which incorporates the irrealis marker \textit{by}. This segment has no separate morpheme status, although it keeps the grammatical requirements of the subjunctive marker \textit{by} from which it derives (see Footnote 13). These requirements are the \textit{l}-form (with a subject in the nominative, see (\ref{ex:wiemer7b})) or the infinitive (with a subject in the dative, see (\ref{ex:wiemer7a})). In the infinitive construction the dative subject is optional; without it the construction implies self-reference to the speaker (or a group including the speaker).\footnote{In narrative contexts, instead of the speaker (narrator), some protagonist may become the vantage point for self-reference. Compare examples (\ref{ex:wiemer2}), (\ref{ex:wiemer20}), (\ref{ex:wiemer21}), (\ref{ex:wiemer24}), (\ref{ex:wiemer27}), (\ref{ex:wiemer30}), (\ref{ex:wiemer31}), (\ref{ex:wiemer41}), (\ref{ex:wiemer56}).} The \textit{l}-form, in turn, can be classified as a morphone, i.e.\ a property of word forms that bears no effect on syntax or semantic interpretation\footnote{The \textit{l}-form is originally an active anteriority participle which has turned into the marker of generalized past tense in North Slavic and western South Slavic. It takes part in subjunctive marking only together with \textit{by}, which is the paradigmatically isolated \textsc{3sg}-form of the former aorist of \textit{byti} `be', originally used as an auxiliary to mark the pluperfect in combination with the \textit{l}-participle. After the aorist and, consequently, the pluperfect were lost, the remnants of \textit{by+l}-participle combinations were reinterpreted as subjunctive (i.e.\ an irrealis mood). In this already intransparent structure \textit{by} still behaved as a Wackernagel clitic, it therefore continued attaching to clause-initial units (WH-words and connectives). With many of them it has coalesced, among others with Russian \textit{kak} `how, that' and other connectives to become complementizers or `optative particles'. \textit{By} and \textit{l}-form need not appear adjacent to each other; the \textit{l}-form requirement of \textit{by} has persisted even after coalescence and the concomitant loss of morpheme status of the latter. For the stages of this process and some diachronic background \mbox{cf.\ \citet{Wiemer2015}}.} \textrm{(\citealt{Aronoff1994}; \citealt{Stump2016}). The segment} \textrm{\textit{by}} \textrm{follows either a WH-word (Russ.} \textrm{\textit{kak}} \textrm{`how') or a particle; among the latter we find focus particles (Russ.} \textrm{\textit{liš'}}, \textrm{\textit{tol'ko}} \textrm{`only') and a particle with concessive function (Russ.} \textrm{\textit{xot'}}). All these combinations also occur without negation and are then used as optative markers (often in exclamations). The units discussed here (with negation) can thus be considered negative-optative markers proper.

An alternative involves connectives that are otherwise prominently employed as complementizers or subordinators of purpose clauses (e.g., Russian \textit{čtoby}). Such connectives thus indicate a connection to negative purpose. In Russian, however, this variant has become obsolete (see \sectref{subsubsec:negpurpose}), while its equivalent in Polish (\textit{żeby}) is used more widely. West Slavic languages also exhibit particles in which \textit{by} has fused with coordinative conjunctions; compare Polish/Czech \textit{aby} (\citealt{lichtenberk1995}: 310 for Czech) as in \REF{ex:wiemer8}.


\ea \textnormal{Czech} \\
\label{ex:wiemer8} 
\gll Aby nám ne-prše-l-o na cest-u. \\
\textsc{con.irr} \textsc{1pl.dat} \textsc{neg}-rain\textsc{[ipfv]-pst-n} on road-\textsc{acc} \\
\glt `[I fear] it might rain while we're on the road.'
\z


In West Slavic, in turn, apprehensional constructions with focus or concessive particles appear to be lacking. The variant with infinitive and dative subject (\ref{ex:wiemer7a}) is altogether lacking (Czech, Slovak) or avoided (Polish), at least in independent clauses (see \sectref{subsec:embedding}). 

In the contemporary languages, the relation to optative and purpose clauses differs. In Russian, constructions of the type (\ref{ex:wiemer7a})--(\ref{ex:wiemer7b}) are compositional \mbox{inasmuch} as the negation is included in the scope of an optative structure (cf.\ \citealt{Dobrushina2016}: 339--342); for non-embedded clauses compare (\ref{ex:wiemer7a})--(\ref{ex:wiemer7b}) with (\ref{ex:wiemer9a})--(\ref{ex:wiemer9b}) and (\ref{ex:wiemer10a})--(\ref{ex:wiemer10b}). However, clause-initial \textit{kak by} (+ \textit{l}-form/infinitive) can induce other illocutions (combined with corresponding intonations) than optatives (or mirative exclamation). First of all, it can be used as a question, with \textit{by} weakening epistemic commitment (see interpretation (i) for (\ref{ex:wiemer10})).

\newpage
\ea \textnormal{Russian}
\label{ex:wiemer9} 
\ea \label{ex:wiemer9a}
\gll Liš' by {\op}nam{\cp} uspe-t'. \\
only \textsc{irr} \textsc{1pl.dat} manage\textunderscore in\textunderscore time\textsc{[pfv]-inf} \\
\ex \label{ex:wiemer9b}
\gll Liš' by my uspe-l-i. \\
only \textsc{irr} \textsc{1pl.nom} manage\textunderscore in\textunderscore time\textsc{[pfv]-pst-pl}\\
\glt `If only we come in time.' / `May we come in time.'
\z
\z



\ea \textnormal{Russian}
\label{ex:wiemer10} 
\ea \label{ex:wiemer10a}
\gll Kak by (nam) uspe-t'? \\
how \textsc{irr} \textsc{1pl.dat} manage\textunderscore in\textunderscore time\textsc{[pfv]-inf} \\
\ex \label{ex:wiemer10b}
\gll Kak by my uspe-l-i? \\
how \textsc{irr} \textsc{1pl.nom} manage\textunderscore in\textunderscore time\textsc{[pfv]-pst-pl}\\
\glt (i) `How might it be possible (for us) to come in time?', \\
(ii) `How can we manage to come in time?', `Can we really (manage to) come in time?'
\z
\z


 In modern West Slavic, constructions like (\ref{ex:wiemer8}) show no obvious relation to manner or exclamative clauses, nor to optative clauses (contrary to Russian \textit{kak by ne}). The relation to purpose is evident only for embedded apprehensional clauses (see \sectref{subsec:embedding}). This is to be expected since, contrary to optatives, purpose clauses usually imply embedding (see \sectref{subsubsec:negpurpose}).

Note that optatives and apprehensionals are polar opposites only insofar as Component 2 of \tabref{tab:componentsslavic} is concerned: wish and apprehension contrast in whether the speaker considers the situation desirable or undesirable (\citealt{Dobrushina2016}: 339--341, \citealt{Wiemer2017}: 298--299), but optatives are indifferent to epistemic attitudes (a wish may be realistic or not), i.e.\ the speaker need not have any specific commitment to the likelihood of the desired situation. On another level, namely in terms of emotional evaluation, wishes are polar to curses, whereas curses and apprehension are, in turn, polar to each other when we consider the relation to the addressee:\footnote{Moreover, curses are always (and wishes can be) uttered in performative speech acts, while apprehension can hardly be imagined as a performative speech act (probably because of its epistemic component).} curses are uttered if the speaker wishes something bad to the addressee, whereas apprehension implies the speaker's concern for the addressee's benefit (see also \sectref{subsec:negimp}).


Regardless of this, in both the apprehensional and the corresponding optative structure all involved elements exist as separate units, but only their combination yields these specific illocutionary functions. The same applies, \textit{mutatis mutandis}, for apprehensional structures and their relation to purpose clauses in West Slavic.


In Balkan Slavic the situation differs, since clauses with initial \textit{da ne} can be variably used either as negative-optatives (${\rightarrow}$ focus on component (a); see (\ref{ex:wiemer11})) or as negative-directives (${\rightarrow}$ focus on component (b); see (\ref{ex:wiemer4}) from Bulgarian and (\ref{ex:wiemer12}) from Macedonian). Apart from their negative-optative function, \textit{da ne}-clauses can also be used as negative-purpose clauses ( \ref{ex:wiemer13}), and this may be taken as an indication that for (\textsc{NegOpt} strategies there is no clear distinction between main and subordinate clause (see further in \sectref{subsec:embedding} and \sectref{subsec:insubordination}). The specific function is usually determined only by the discourse context. However, the negative-purpose function results from a transparent compositional reading of \textit{da} as a general irrealis morpheme (with purpose being one possibility) and the negation.


\ea
\label{ex:wiemer11} 
\langinfo{Macedonian}{}{\citealt[68]{MitkovskaEtAl2017}}\\
\gll Aj dosta, Boško, \textbf{da} \textbf{ne} me pobara tatko $\ldots$\\
\textsc{ptc} enough \textsc{pn} \textsc{irr} \textsc{neg} \textsc{1sg.acc} call\textsc{[pfv].prs.3sg} dad \\
\glt `Boško, I must go, \textbf{lest} my father call for me $\ldots$' \\
$\rightarrow$ (i) negative-optative; (ii) \sout{preventive}
\z

\ea
\label{ex:wiemer12} 
\langinfo{Macedonian}{}{\citealt[74]{MitkovskaEtAl2017}} \\
\gll \textbf{Da} \textbf{ne} nastine-š! \\
\textsc{irr} \textsc{neg} get\textunderscore cold\textsc{[pfv]-prs.2sg}\\
\glt (i) `May you not get a cold.' (ii) `(Be careful!) Don't catch a cold!' \\
$\rightarrow$ (i) negative-optative (ii) preventive [preferred]
\z

\ea
\label{ex:wiemer13} 
\langinfo{Macedonian}{}{\citealt[66]{MitkovskaEtAl2017}} \\
\gll Duva-j mil-o, \textbf{da} \textbf{ne} se popari-š. \\
blow\textsc{[ipfv]-imp.sg} dear-\textsc{voc.sg.m} \textsc{irr} \textsc{neg} \textsc{refl} scald\textsc{[pfv]-prs.2sg}\\
\glt `Run, dear, (i) or you might get scаlded; (ii) so as not to be scalded.' \\
(i) negative-optative; (ii) negative-purpose ($>$ preventive) \\
\z


Negative purpose and preventive readings can become indistinguishable (see (\ref{ex:wiemer13}) and also (\ref{ex:wiemer37}) below), and this applies to pertinent constructions in West Slavic as well (see, for instance, the Polish examples (\ref{ex:wiemer21}) and (\ref{ex:wiemer22}) below). The reasons for this are of a conceptual nature: (negative) purpose and prevention both presuppose control, while (negative) optatives do not (see \sectref{sec:conprem}). Since, in addition, South Slavic constructions based on \textit{da ne} (despite their different opaque vs.\ transparent structure) do not discriminate between negative purpose and negative optative function (and they are anyway very vague in their meaning range; see \sectref{subsec:lossofcomp}), they can in principle cover also preventive functions. Whether diachronically this is to be seen as an extension from negative purpose, is an open issue (see \sectref{subsubsec:balkanslavic}).


The ``flexible'' behavior of \textit{da ne} is, to a certain extent, inherited from the vague functional values of \textit{da}, which from the earliest attested times (in Old Church Slavonic) was a general irrealis-marker; optative use was among its salient uses (\citeauthor{Večerka1993} \citeyear[79]{Večerka1993}, \citeyear[65]{Večerka1996}, \citealt{Grković-Major2020}). Balkan Slavic continues this situation (see (\ref{ex:wiemer14})--(\ref{ex:wiemer16})), but, in addition, \textit{da} has become part of a cluster of verb-adjacent proclitics and it constrains the admissible range of TAM forms. Depending on the construction and the context of speech, the illocutionary function of \textit{da} is vague, and its syntactic status in clause combining is very versatile (\citeauthor{Wiemer2017} \citeyear[300--305, 325--330]{Wiemer2017}, \citeyear[295--306]{Wiemer2018}, \citeyear[58--80]{Wiemer2021}). This carries over to combinations with negation (\textit{ne}); among other things, \textit{da ne} can be used as an adverbial subordinator (conjunction) with purpose clauses. \textit{Da ne} shows semantically transparent usage types, in particular of [optative+negation $>$ apprehensional], although in other environments (e.g.\ as an epistemic modifier) \textit{da ne} loses practically all the aforementioned properties, and this correlates with the loss of morphological autonomy and semantic transparency (see \sectref{subsec:lossofcomp}). The relation between apprehensional and optative function becomes obvious when we look at equivalent utterances without negation, for instance\footnote{See \citet{Tomić2012} for a comprehensive survey of imperatives and ``bare subjunctives''; cf.\ also \citet{Mitkovska2015} and \citet{MitkovskaEtAl2017}, with further references concerning imperative-like uses.} (see also (\ref{ex:wiemer37}) vs.\ (\ref{ex:wiemer38}) in \sectref{subsec:lossofcomp}):


\ea
\label{ex:wiemer14} 
\langinfo{Macedonian}{}{\citealt[371]{Tomić2012}}\\
\gll Da gi prečeka-te! \\
\textsc{irr} \textsc{3pl.acc} await\textsc{[pfv]-prs.2pl} \\
\glt `(You should) Welcome them!'
\z


\ea
\label{ex:wiemer15} 
\langinfo{Macedonian}{}{\citealt[371]{Tomić2012}}\\
\gll Golem da raste-š! \\
big\textsc{.sg.m} \textsc{irr} grow\textsc{[ipfv]-prs.2sg} \\
\glt `May you grow big!'
\z

\ea
\label{ex:wiemer16} 
\langinfo{Bulgarian}{}{\citealt[2]{Mitkovska2015}} \\
\gll Da băd-eš zdrav i vesel! \\
\textsc{irr} be\textsc{[pfv]-prs.2sg} healthy\textsc{.sg.m} and cheerful\textsc{.sg.m} \\
\glt `May you be healthy and cheerful!' \\

\z


\noindent These utterances have an optative (\ref{ex:wiemer15})--(\ref{ex:wiemer16}) or hortative (\ref{ex:wiemer14}) value. In general, Balkan Slavic \textit{da}+indicative can be used as an equivalent of the imperative in its different context-conditioned functions (see Footnote 15).


Basically the same comments are justified for Bulgarian \textit{da ne bi} (etymologically \textsc{irr+neg+sbjv}); compare (\ref{ex:wiemer2}) and the negative purpose clause in (\ref{ex:wiemer17}):


\ea
\label{ex:wiemer17} 
\langinfo{Bulgarian}{}{\citealt[64]{MitkovskaEtAl2017}} \\
\gll \textnormal{($\ldots$)} tutaksi ugasi fener-a \textbf{da} \textbf{ne} \textbf{bi} njakoj ot reka-ta da go zabeleži.	\\
{} immediately extinguish\textsc{[pfv].aor.3sg} flashlight-\textsc{def.do} \textsc{irr} \textsc{neg} \textsc{sbjv} somebody from river-\textsc{def} \textsc{irr} \textsc{3m.acc} notice\textsc{[pfv].prs.3sg} \\
\glt `He immediately turned the flashlight off \textbf{so that} no one could notice him from the river.' \\
\z


 However, \textit{da ne bi} differs from \textit{da ne} in that it normally requires another \textit{da} proclitic to the verb (see further \sectref{subsec:lossofcomp}). In addition, it includes another irrealis marker \textit{bi} deriving from an older Slavic subjunctive (or conditional). The form \textit{bi} is fossilized (originally \textsc{3sg}) and does not show agreement, in contrast to the subjunctive marker proper\footnote{Both \textit{bi} and \textit{by} (used in North Slavic) relate to copular \textit{byti} `be', but \textit{bi} does not derive from a former aorist (see Footnote 13). Moreover, in contrast to the North Slavic \textit{by+l} form, the subjunctive marked with the \textit{bi+l} form has become marginal particularly in Balkan Slavic, being more and more superseded by \textit{da}+\textsc{(pfv)present}. Combinations of \textit{da+bi} have found lexicalized niches in formulaic expressions (e.g., curses).} \citep[147]{Ivanova2014}. \textit{Da ne bi} acquires apprehensional meaning preferably in embedded clauses \citep[143]{Ivanova2014}.
\noindent However, \textit{da ne bi} differs from \textit{da ne} in that it normally requires another \textit{da} proclitic to the verb (see further \sectref{subsec:lossofcomp}). In addition, it includes another irrealis marker \textit{bi} deriving from an older Slavic subjunctive (or conditional). \textit{Da ne bi} acquires apprehensional meaning preferably in embedded clauses \citep[143]{Ivanova2014}.
\noindent However, \textit{da ne bi} differs from \textit{da ne} in that it normally requires another \textit{da} proclitic to the verb (see further \sectref{subsec:lossofcomp}). In addition, it includes another irrealis marker \textit{bi} deriving from an older Slavic subjunctive (or conditional). The form \textit{bi} is fossilized (originally \textsc{3sg}) and does not show agreement, in contrast to the subjunctive marker proper\footnote{Both \textit{bi} and \textit{by} (used in North Slavic) relate to copular \textit{byti} `be', but \textit{bi} does not derive from a former aorist (see Footnote 13). Moreover, in contrast to the North Slavic \textit{by} + \textit{l}-form, the subjunctive marked with the \textit{bi}+ \textit{l}-form has become marginal particularly in Balkan Slavic, being more and more superseded by \textit{da}+\textsc{(pfv)present}. Combinations of \textit{da+bi} have found lexicalized niches in formulaic expressions (e.g., curses).} \citep[147]{Ivanova2014}. \textit{Da ne bi} acquires apprehensional meaning preferably in embedded clauses \citep[143]{Ivanova2014}.




Both Macedonian \textit{da ne} and Bulgarian \textit{da ne bi} allow for `in-case' readings (see (\ref{ex:wiemer18})); these are excluded for \textsc{NegOpt} in North Slavic:


\ea
\label{ex:wiemer18} 
\langinfo{Macedonian}{}{\citealt[67]{MitkovskaEtAl2017}} \\
\gll Ќe mu se java-m na Marko \textbf{da} \textbf{ne} ne zna-e za sostanok-ot. \\
\textsc{fut} \textsc{3m.dat} \textsc{refl} call\textsc{[pfv]-prs.1sg} on \textsc{pn} \textsc{irr} \textsc{neg} \textsc{neg} know\textsc{[ipfv]-prs.3sg} for meeting\textsc{.sg.m-def.sg.m} \\
\glt `I'll call Marko, \textbf{in case} he doesn't know about the meeting.' \\
\z


 In sum, Balkan Slavic clause-initial particles used in \textsc{NegOpt} show a wider range of uses than their North Slavic counterparts. Among other things, they are employed as rhetoric discourse markers (with mirative, epistemic and other functions). Also from a morphological point of view the \textsc{NegOpt} strategies are much less distinguishable from \textsc{NegDir} strategies (preventives) than they are in North Slavic (see \sectref{subsec:lossofcomp}).


Moreover, in Russian the structure of \textsc{NegOpt} varies depending on whether the clausal subject differs from the speaker or whether both coincide. If it differs, the sentence has a morphologically finite predicate (see (\ref{ex:wiemer1}), (\ref{ex:wiemer7b})); if it coincides, the predicate is predominantly an infinitive with an (optional) dative subject (see (\ref{ex:wiemer7a})). In West Slavic the infinitival construction is avoided (in main clauses), while in Balkan Slavic the infinitive is lost altogether,\footnote{Remnants of a shortened infinitive in Bulgarian are rare and irrelevant for apprehensional constructions.} so that surface distinctions between the two constellations are precluded. Compare (\ref{ex:wiemer2}) and (\ref{ex:wiemer11}) (no coincidence) with (\ref{ex:wiemer19}) (coincidence):

\ea
\label{ex:wiemer19} 
\langinfo{Macedonian}{}{\citealt[64]{MitkovskaEtAl2017}} \\
\gll Najmnogu se plaš-ev \textbf{da} \textbf{ne} ja razočara-m. \\
{most of all} \textsc{refl} fear\textsc{[ipfv]-impf.1sg} \textsc{irr} \textsc{neg} \textsc{3f.acc} disappoint\textsc{[pfv]-prs.1sg} \\
\glt `Most of all I was afraid not to let her down.' (lit.\ `$\ldots$ lest I disappoint her.') \\
\z



\subsection{Embedding}
\label{subsec:embedding}

\textsc{NegOpt} can occur as a clausal complement of verbs denoting fear (or related notions), of related nouns (for the latter see (\ref{ex:wiemer25})--(\ref{ex:wiemer26})) or of predicates of adjectival origin (e.g., Russian \textit{bojazno} `(it is) fearful, frightening').\footnote{For lists of such verbal and nominal predicates in contemporary Russian cf.\ \citet[56]{Nilsson2012}, \citet[338]{Dobrushina2016}.} Embedding implies that clause-initial units like Russian \textit{kak by $\ldots$ ne}, Bulgarian \textit{da ne bi}, Polish \textit{żeby $\ldots$  nie} function as complementizers (see \sectref{subsec:tempref} and \sectref{subsec:insubordination}). Not all clause-initial connectives with apprehensional function occur as complementizers; for instance, Russian \textit{liš' by ne}, \textit{tol'ko by ne} and \textit{xot' by ne} are not attested in the RNC as complementizers, while \textit{kak by $\ldots$ ne} is more frequent in clauses that can be qualified as complements than in independent clauses (\citealt[335 and the Table on p. 330]{Dobrushina2016}). However, \textit{kak by $\ldots$  ne} still reveals more properties of autonomy than canonical complementizers (\citealt[345-346]{Dobrushina2016}); this correlates with its probably very recent complementizer status (\sectref{subsec:insubordination}).

Embedding turns the distinction of referential identity vs.\ difference between speaker and clausal subject into a distinction of different-subject (DS) vs.\ same-subject (SS) between the subjects of the main and the subordinate clause. Example (\ref{ex:wiemer20}) illustrates DS.


\ea
\label{ex:wiemer20} 
\langinfo{Bulgarian}{}{\citealt[150]{Ivanova2014}} \\
\gll Toj \textit{se} \textit{panik'osa} \textbf{da} \textbf{ne} \textbf{bi} gen. Borisov da potărsi văzmezdie za minalo-to. \\
\textsc{3m.nom} \textsc{refl} panic\textsc{[pfv].aor.3sg} \textsc{irr} \textsc{neg} \textsc{sbjv} general \textsc{pn} \textsc{irr} want\textsc{[pfv].prs.3sg} revenge for past\textsc{[n]-def.n} \\
\glt `He \textit{panicked} \textbf{lest} general Borisov wants to take revenge for the past.' \\
\z


 For embedded \textsc{NegOpt} clauses, contrasts of morphological finiteness, as encountered in Russian, correspond to a distinction of balanced (for DS) vs.\ deranked (for SS) clause-combining (following \citealt{Cristofaro2003}). With SS, infinitival, i.e.\ de\-ranked, complements are practically the only option in Russian (see \sectref{subsec:tempref}). However, as a complementizer of apprehensive clauses \textit{kak by ne} is overall more frequent with the \textit{l}-form (\citealt[338]{Dobrushina2016}, based on the RNC), and this pattern is strongly associated with DS (see \sectref{subsec:tempref}).

As for West Slavic, Polish allows for infinitival complements in case of SS (see (\ref{ex:wiemer21})), although they are not very frequent. Moreover, the negative-purpose and preventive readings are often hardly distinguishable (see (\ref{ex:wiemer22})):


\ea \textnormal{Polish}\\
\label{ex:wiemer21} 
\gll Wymy-ł-a puszk-ę piask-iem i napełni-ł-a ją wod-ą, uważa-jąc, \textbf{żeby} \textbf{jej} \textbf{nie} \textbf{zmąci-ć.} \\
wash\textsc{[pfv]-pst-3.sg.f} can\textsc{[f]-acc.sg} sand\textsc{[m]-ins.sg} and fill\textsc{[pfv]-pst-3.sg.f} \textsc{3.acc.sg.f} water\textsc{[f]-ins.sg} take\textunderscore care\textsc{[ipfv]-cvb} \textsc{con.irr} \textsc{3f.gen} \textsc{neg} disturb\textsc{[pfv]-inf} \\
\glt `She washed the can with sand and filled it with water, \textit{taking care} \textbf{not to disturb her}.' (`$\ldots$  lest she disturb her.')\\
 {[PNC; I.\ Jurgielewiczowa: \textit{Ten obcy}. 1990 [1961]]}
\z



\begin{exe} \ex\textnormal{Polish}
\label{ex:wiemer22} \\
\gll Pilnowa-ł-em się, \textbf{żeby} \textbf{nieopatrznie} \textbf{cz-ego{\' s}} \textbf{nie} \textbf{chlapną-ć} o jej męż-u. \\
beware\textsc{[ipfv]-pst-1sg.m} \textsc{refl} \textsc{con.irr} unattentively something-\textsc{gen} \textsc{neg} blurt\textunderscore out\textsc{[pfv]-inf} about her husband-\textsc{loc} \\
\glt `\textit{I was careful} \textbf{not to blurt out anything inadvertently} about her husband.' \\
(i) `$\ldots$  in order to not blurt out $\ldots$'; (ii) `$\ldots$ lest I blurt out $\ldots$'\\ {[PNC; J.\ Grzegorczyk: \textit{Chaszcze.} 2009.]}
\end{exe}

 Czech and Slovak, contrary to Polish, do not allow the irrealis connective (\textit{aby}; see (\ref{ex:wiemer8})) to introduce infinitival complements (M. Giger, p.c.); the use of infinitival clauses in these two languages is generally very restricted. Thus, the only opposition is between \textit{aby}, which must combine with (so-called expletive) negation and the \textit{l}-form, and a default complementizer (Czech/Slovak \textit{že}), which goes without expletive negation and does not restrict the TAM-choice in the verb. These options exist in Russian and Polish as well, but their choice does not affect the SS-DS distinction. It is however indicative of the balanced-deranked distinction, provided grammatical distribution is taken into account: the option with irrealis connective and negation (\textit{aby $\ldots$ ne}) restricts the choice of the verb form (for reasons given in Footnote 13), while the default complementizer does not (see \sectref{subsec:tempref}). However, these restrictions only follow from the general collocational properties of clause-initial connectives with incorporated irrealis marker (\textit{by}), which are the same in independent clauses.


Analogous remarks apply to Balkan Slavic embedded \textsc{NegOpt} clauses. Al\-though a finite-infinitive contrast is not available, deranking is based on grammatical distribution: the predicate of an embedded apprehensional clause very rarely occurs in the aorist; example (\ref{ex:wiemer23}) is constructed, since even an internet search does not provide good examples. Instead, the predicate of an embedded apprehensional clause predominantly takes a present tense form; more occasionally it occurs in the perfect.\footnote{A future marker (Bulg.\ \textit{šte}, Mac.\ \textit{\textsc{\'k}e}) is excluded, since it never co-occurs with \textit{da} in the same cluster.} See (\ref{ex:wiemer24}), which shows simultaneous reference, whereas (\ref{ex:wiemer25}) illustrates anterior reference:


\ea
\label{ex:wiemer23} 
\langinfo{Macedonian}{}{constructed, by courtesy of E.\ Bužarovska and L.\ Mitkovska} \\
\gll Se plaš-am \textbf{da} \textbf{ne} zapadn-aa vo nevolja. \\
\textsc{refl} be\textunderscore afraid\textsc{[ipfv]-prs.1sg} \textsc{irr} \textsc{neg} fall\textsc{[pfv]-aor.3pl} into trouble \\
\glt `I fear \textbf{that} they have got into trouble.' \\
\z


\ea
\label{ex:wiemer24} 
\langinfo{Macedonian}{}{\citealt[67]{MitkovskaEtAl2017}} \\
\gll Tina molče-še, isplaše-n-a \textbf{da} \textbf{ne} ima i taa takov virus. \\
\textsc{pn} be\textunderscore quiet\textsc{[ipfv]-impf.3sg} frighten\textsc{[pfv]-pp-sg.f} \textsc{irr} \textsc{neg} have\textsc{[ipfv].prs.3sg} also \textsc{3f.nom} such virus\textsc{.sg.m} \\
\glt `Tina kept quiet, fearing \textbf{that} she \textbf{might} have the same virus.' \\
\z


\ea
\label{ex:wiemer25} 
\langinfo{Bulgarian}{}{\citealt[67]{MitkovskaEtAl2017}} \\
\gll Strax me e \textbf{da} \textbf{ne} \textbf{bi} da ne e doš-l-a. \\
fear \textsc{1sg.acc} be.\textsc{prs.3sg} \textsc{irr} \textsc{neg} \textsc{sbjv} \textsc{irr} \textsc{neg} be\textsc{.prs.3sg} arrive\textsc{[pfv]-lf-sg.f} \\
\glt `I fear \textbf{that} she \textbf{might} not have come.'
\z

 Deranking even in this extended sense does not seem to apply in the western part of South Slavic (Serbian-Croatian, Slovene), as embedded apprehensional clauses allow for future and conditional mood (see \sectref{subsec:tempref}). This appears to follow from the loss of irrealis restrictions otherwise imposed by \textit{da}.

\subsection{Temporal reference}
\label{subsec:tempref}

Since in North Slavic clause-initial connectives with an incorporated irrealis marker (\textit{by}) restrict the form of the verb (\textit{l}-participle or, in Russian and Polish, infinitive), clauses containing them cannot distinguish temporal reference, either to speaker deixis (in independent use) or to the reference interval set by a matrix predicate for an embedded apprehensional clause. Nonetheless, if \textit{by} participates in the formation of \textsc{NegOpt} strategies, these show a strong tendency toward posterior reference (to either a deictic or a syntactically anchored reference interval). Russian (independent or dependent) infinitive clauses only allow for a posterior orientation (see (\ref{ex:wiemer9a}) and (\ref{ex:wiemer26b})--(\ref{ex:wiemer26bsecond})), but even with the \textit{l}-form, posterior orientation seems to practically be the  exclusive option. Compare the following textbook examples:

\newpage
\ea \textnormal{Russian} \\
\label{ex:wiemer26} 
\ea \label{ex:wiemer26a} \textnormal{DS: biclausal with different subjects $\rightarrow$ balanced (symmetric):} \\
\gll Ja boj-u-s', \textbf{kak} \textbf{by} on-a \textbf{ne} prostudi-\textbf{l}-a-s'. \\
\textsc{1sg.nom} fear\textsc{[ipfv]-prs.1sg-refl} how \textsc{irr} 3-\textsc{f.nom} \textsc{neg} catch\textunderscore cold\textsc{[pfv]-pst-sg.f-refl} \\
\glt `I fear that she (i) ?has caught a cold. (ii) `$\ldots$ will catch a cold.'
\ex \label{ex:wiemer26b} \textnormal{SS: biclausal with identical subjects $\rightarrow$ deranked (asymmetric):} \\
\gll Ja boj-u-s', \textbf{kak} \textbf{by} \textbf{ne} prostudi-t'-sja. \\
\textsc{1sg.nom} fear\textsc{[ipfv]-prs.1sg-refl} how \textsc{irr} \textsc{neg} catch\textunderscore cold\textsc{[pfv]-inf-refl} \\
\ex \label{ex:wiemer26bsecond}
\gll = Ja boj-u-s' prostudi-t'-sja. \\
{} \textsc{1sg.nom} fear\textsc{[ipfv]-prs.1sg-refl} catch\textunderscore cold\textsc{[pfv]-inf-refl} \\
\glt `I fear to catch a cold / that I will catch a cold.'
\z
\z


\noindent While the embedded infinitive clause in (\ref{ex:wiemer26b}) and (\ref{ex:wiemer26bsecond}) induces a posterior reading, the finite embedded clause in (\ref{ex:wiemer26a}) might in principle allow for an anterior reading (as indicated in (i)). However, native speakers are highly reluctant to accept this interpretation (\citealt[282]{Dobrushina2016} considers it inappropriate), and I have been unable to find convincing examples in the RNC (nor does \citealt{Nilsson2012} adduce any). Only theoretically is there an alternative to (\ref{ex:wiemer26b}) with a finite predicate:


\ea 
\label{ex:wiemer26c} \textnormal{Russian} \\
\gll Ja boj-u-s', \textbf{kak} \textbf{by} ja \textbf{ne} prostudi-\textbf{l}-\emptyset-sja. \\
\textsc{1sg.nom} fear\textsc{[ipfv]-prs.1sg-refl} how \textsc{irr} \textsc{1sg.nom} \textsc{neg} catch\textunderscore cold\textsc{[pfv]-pst-sg.m-refl} \\
\glt `I fear that I will catch a cold / ?$\ldots$ that I have caught a cold.'
\z

\noindent Since informants judge such an option as artificial, speculations about its possible interpretation are senseless. We may thus presume that, in Russian, the practically complementary distribution of (DS ${\times}$ finite) vs.\ (SS ${\times}$ non-finite) predicates is not accompanied by variability of temporal orientation (anterior -- simultaneous -- posterior), i.e.\ the structure \textit{kak} \textit{by $\ldots$ ne} is strongly associated with posterior reference in both dependent and independent clauses.


The situation differs in South Slavic, but differently for different subareas. In the (north)western part, i.e.\ in Slovene and (standard) Serbian-Croatian, \textit{verba timendi} combine with \textit{da} as default complementizer (`that'), which has basically lost its irrealis semantics and does not function as a verb-adjacent proclitic.\footnote{For this process and its inner-South Slavic cline cf.\ \citet{Sedláček1970}; \citet{Grković-Major2004}; \citet{Grković-Major2020}. For an account of Slovene cf.\ \citet[98--103]{Žele2017}. All examples and claims on Slovene in the remainder of this subsection are taken from there.} Concomitantly, in Slovene verbs like \textit{bati se} `be afraid' do not restrict the TAM-form of the predicate in their dependent clause:


\ea
\label{ex:wiemer27} \textnormal{Slovene} \\
\gll Politik \textbf{se} \textbf{boji,} $\ldots$ \\
politician\textsc{.nom.sg} \textsc{refl} fear\textsc{[ipfv].prs.3sg} \\
\ea \label{ex:wiemer27a} \textnormal{Past (with \textit{l}-form)} \\
\gll da \textbf{so} volilc-i \textbf{izvoli-l-i} nasprotnik-a. \\
\textsc{comp} \textsc{be.prs.3pl} voter-\textsc{nom.pl} vote\textsc{[pfv]-lf-pl} opponent-\textsc{acc} \\
\ex \label{ex:wiemer27b} \textnormal{Present} \\
\gll da \textbf{volij-o} nasprotnik-a. \\
\textsc{comp} vote\textsc{[ipfv]-prs.3pl} opponent-\textsc{acc} \\
\ex \label{ex:wiemer27c} \textnormal{Future} \\
\gll da \textbf{bo-do} \textbf{izvoli-l-i} nasprotnik-a. \\
\textsc{comp} \textsc{fut-3pl} vote\textsc{[pfv]-lf-pl} opponent-\textsc{acc} \\
\glt `The politician is afraid $\ldots$' \\
(a) `that the voters have voted for the opponent.' \\
(b) `that they vote for the opponent.' \\
(c)	`that they will vote for the opponent.'
\z
\z

\noindent Forms of the indicative future (\ref{ex:wiemer27c}) seem to predominate. Moreover, as (\ref{ex:wiemer27a}) to (\ref{ex:wiemer27c}) show, \textit{da} tends to go without negation (\textit{ne}). Negation is not excluded, but the factors that condition its interpretation as ``ordinary'' sentence negation, as in (\ref{ex:wiemer28}), as opposed to expletive negation, have not been clarified yet. In either case, \textit{da} can occur with the future (\ref{ex:wiemer28}) or the subjunctive (\ref{ex:wiemer29a}) marker. An infinitival complement is possible, in principle, in SS clause combining (\ref{ex:wiemer29b}). However, in cases of SS, a dependent apprehensional clause tends to be finite (other than in Russian or Polish, but as in Czech and Slovak); compare (\ref{ex:wiemer26c}) with (\ref{ex:wiemer28}).


\ea 
\label{ex:wiemer28} \textnormal{Slovene} \\
\gll Kakšn-a suš-a, boji-m se, da ne bo dolgo deževa-l-o. \\
which-\textsc{nom.sg.f} drought\textsc{[f]-nom.sg} fear\textsc{[ipfv]-prs.1sg} \textsc{refl} \textsc{irr} \textsc{neg} \textsc{fut.3sg} long\textunderscore time rain\textsc{[ipfv]-lf-n} \\
\glt `What a drought. I'm afraid it won't rain for long.'
\z


\ea
\label{ex:wiemer29} \textnormal{Slovene} \\
\ea \label{ex:wiemer29a} \gll Boji-m se, da (ne) bi pozabi-l-\emptyset {} denarnic-o. \\
fear\textsc{[ipfv]-prs.1sg} \textsc{refl} \textsc{comp} \textsc{neg} \textsc{sbjv} forget\textsc{[pfv]-lf-sg.m} purse-\textsc{acc} \\
\glt `I fear that I might / will forget my purse.' \\
\ex \label{ex:wiemer29b}
\gll Boji-m se pozabi-ti denarnic-o. \\
fear\textsc{[ipfv]-prs.1sg} \textsc{refl} forget\textsc{[pfv]-inf} purse-\textsc{acc} \\
\glt `I fear to forget my purse.'
\z
\z



 The situation in Serbian-Croatian is similar, however here perfective present tense with negation appears to be the preferred option in dependent clauses (\ref{ex:wiemer30a}). Another option is modal auxiliaries in the subjunctive (\ref{ex:wiemer30c}); compare the examples in (\ref{ex:wiemer30}).



\ea 
\label{ex:wiemer30}
\langinfo{Serbian}{}{\citealt[163]{Alanović2018}} \\
\gll Boji se, $\ldots$ \\
fear\textsc{[ipfv]-prs.3sg} \textsc{refl} \\
\ea \label{ex:wiemer30a} \textnormal{Perfective present ($+$ negation)} \\
\gll da ne zakasni. \\
\textsc{irr} \textsc{neg} be\textunderscore late\textsc{[pfv].prs.3sg} \\
\ex \label{ex:wiemer30b} \textnormal{Perfective future} \\
\gll da će zakasni-ti. \\
\textsc{irr} \textsc{fut.3sg} be\textunderscore late\textsc{[pfv]-inf} \\
\ex \label{ex:wiemer30c} \textnormal{Conditional with possibility modal auxiliary} \\
\gll da bi moga-o-\emptyset {} zakasni-ti. \\
\textsc{irr} \textsc{sbjv} can-\textsc{lf-sg.m} be\textunderscore late\textsc{[pfv]-inf} \\
\glt `He is afraid $\ldots$' \\
(a) `to arrive late' [lit.\ `lest he arrive late']. \\
(b) `that he will arrive late.' \\
(c)	`that he might arrive late.'
\z
\z

 In Balkan Slavic, since the relevant clause-initial particles do not entail any restrictions of temporal orientation, embedded clauses may also show simultaneous or anterior time-reference; compare (\ref{ex:wiemer23})--(\ref{ex:wiemer25}) with (\ref{ex:wiemer31}):


\ea
\label{ex:wiemer31} 
\langinfo{Bulgarian}{}{\citealt[150]{Ivanova2014}} \\
\gll Vrabče-ta-ta pomisli-l-i \textbf{da} \textbf{ne} \textbf{bi} da e izbuxna-l-a vojna-ta, \textnormal{($\ldots$)}\\
sparrow-\textsc{pl-def.pl} think\textsc{[pfv]-lf-pl} \textsc{irr} \textsc{neg} \textsc{sbjv} \textsc{irr} be.\textsc{prs.3sg} break\textunderscore out\textsc{[pfv]-lf-sg.f} war\textsc{[f]-def.sg.f} \\
\glt `The sparrows thought [> were afraid] \textbf{that} the war \textbf{might} have broken out (but nowhere could any rumbling be heard).' \\
\z




\subsection{Loss of compositionality}
\label{subsec:lossofcomp}

All clause-initial connectives relevant for \textsc{NegOpt} show a high degree of coalescence (which has diminished or eliminated morphological transparency) among their components. In South Slavic, these components, first of all, cannot be separated by other (even clitic) units; in North Slavic, \textit{by} has coalesced with the preceding component, but the negation can be separated from it. As for semantic transparency (vs.\ opacity), the matter is more complicated.


Macedonian \textit{da ne} (\textsc{irr+neg}) has lost segmentability and semantic transparency. This regards its use both in apprehensional main and dependent clauses as well as its use as epistemic modifier (employed, for instance, in deliberative questions; see (\ref{ex:wiemer34})). The evidence is, first, that in these functions \textit{da ne} is used non-redundantly with sentence negation (\citealt{Bužarovska2017}, \citealt{MitkovskaEtAl2017}: 67, 75 et passim). See (\ref{ex:wiemer18}) for an `in-case'-reading and (\ref{ex:wiemer32}), which can acquire an apprehensional function:


\ea \textnormal{Macedonian} \\
\label{ex:wiemer32} 
\gll A kade e Marko? \textbf{Da} \textbf{ne} ne dojd-e na vreme? \\
and	where be.\textsc{prs.3sg} \textsc{pn} \textsc{irr} \textsc{neg} \textsc{neg} come\textsc{[pfv]-prs.3sg} on time \\
\glt `Where is Marko? (I'm afraid) maybe he won't come in time.'
\z


 Second, in these functions \textit{da ne} never attracts stress; see the transparent use in (\ref{ex:wiemer36}). Third, the requirement to occur adjacent to the verb (as part of a proclitic cluster) gets lost; see (\ref{ex:wiemer33}) and (\ref{ex:wiemer34}):


\ea 
\label{ex:wiemer33} \textnormal{Macedonian} \\
\gll Se plaš-am \textbf{da} \textbf{ne} \textit{slučajno} dozna-at za toa. \\
\textsc{refl} be\textunderscore afraid\textsc{[ipfv]-prs.1sg} \textsc{irr} \textsc{neg} accidentally learn\textsc{[pfv]-prs.3pl} about this \\
\glt `I'm afraid they might \textit{accidentally} find out about it.'
\z



\noindent Fourth, TAM constraints get lost as well, i.e.\ the aorist is possible (see (\ref{ex:wiemer34})), which, as a ``confirmative tense" \citep{Friedman2000}, is otherwise incompatible with the irrealis function of \textit{da}. As stated in \sectref{subsec:embedding}, in complementizer function this happens very rarely.

It is often difficult to distinguish \textit{da ne} as an epistemic modifier, as in (\ref{ex:wiemer34}) from its use as a marker of apprehension (related to optative or purpose clauses). In fact, \textit{da ne} (in both Macedonian and Bulgarian) can be used as a general epistemic downtoner void of emotional involvement; its apprehensional function can only be inferred from this general function. Compare examples  (\ref{ex:wiemer35}) to (\ref{ex:wiemer38}); stress is indicated by bold capital letters:\footnote{For a detailed treatment cf.\ \citeauthor{Mitkovska2015} (\citeyear[37 et passim]{Mitkovska2015}, \citeyear{MitkovskaEtAl2017}), \citet[67--70]{MitkovskaEtAl2017}.}



\ea 
\textnormal{Macedonian} \\
\ea \label{ex:wiemer34} \textnormal{Epistemic modifier: no apprehension} \\
\gll Da ne \textbf{včera} ja skrši vazn-a-ta? \\
\textsc{irr} \textsc{neg} yesterday \textsc{3f.acc} break\textsc{.aor.2sg} vase[\textsc{f}]-\textsc{sg-def.sg.f} \\
\glt `Did you by any chance break the vase \textbf{yesterday}?' \\
\ex \label{ex:wiemer35} \textnormal{Epistemic modifier: apprehension}\\
\gll Da ne {\op}si{\cp} \textbf{O}di-š? \\
\textsc{irr} \textsc{neg} \textsc{refl.dat} go\textunderscore away\textsc{[ipfv]-prs.2sg} \\
\glt `(I'm afraid:) Are you really leaving?' \\
alternatively, deliberative question: `You are not leaving, are you?' \\
\ex \label{ex:wiemer36} \textnormal{Transparent collocation: prohibitive} \\
\gll Da \textbf{NE} odi-š! \\
\textsc{irr} \textsc{neg} go\textunderscore away\textsc{[ipfv]-prs.2sg} \\
\glt `(I order you:) Don't go!' \\
\ex \label{ex:wiemer37} \textnormal{Opaque: negative optative/negative purpose > preventive} \\
\gll Pobrza-j-te! Da ne go ispušti-te avion-ot! \\
hurry\textsc{[pfv]-imp-pl} \textsc{irr} \textsc{neg} \textsc{3m.acc} miss\textsc{[pfv]-prs.2pl} airplane\textsc{.sg.m.-def.sg.m} \\
\glt `Hurry up! Don't miss the plane!' ($\approx$ `Otherwise you'll miss your plane.') 

\ex \label{ex:wiemer38} \textnormal{Optative} \\
\gll Ubavo da si pomine-š na odmor! \\
nicely \textsc{irr} \textsc{refl.dat} pass\textsc{[pfv]-prs.2sg} on vacation \\
\glt `Have a nice vacation!' (more literally `May you nicely pass for a vacation!')
\z
\z


\noindent In (\ref{ex:wiemer34}) and (\ref{ex:wiemer35}), using \textit{da ne}, the speaker utters an assumption. Likewise, (\ref{ex:wiemer36}) does not express apprehension, whereas (\ref{ex:wiemer37}) does. That is, a deliberative question can yield apprehension only as a side-effect, given a suitable communicative purpose and corresponding intonation.\footnote{Cf.\ \citet{Bužarovska2017} and \citet[68--69]{MitkovskaEtAl2017}, who call this type of question \textit{assumptive}: the speaker wants to gain confirmation for their assumption that the proposition is true.} By contrast, (\ref{ex:wiemer37}) is the negative equivalent of an optative sentence (as in (\ref{ex:wiemer38})), and it can acquire a preventive interpretation. In none of these instances can \textit{da ne} be stressed, contrary to (\ref{ex:wiemer36}), where \textit{da+ne} is a transparent collocation: the negation can be stressed, but the utterance will be understood only as a prohibition. Notably, prohibition is associated rather with imperfective aspect, while a preventive interpretation of a \textsc{NegOpt} construction (as in (\ref{ex:wiemer37})) is rather associated with perfective aspect. This is a tendency, not a rule, but it corresponds to the more clear-cut picture in North Slavic (see \sectref{subsec:negimp}--\sectref{subsec:areal}).


As for Bulgarian \textit{da ne bi} (etymologically \textsc{irr+neg+sbjv}), this unit cannot be split internally, but as a whole it can be separated from the verb by other constituents (see (\ref{ex:wiemer2}), (\ref{ex:wiemer17})). It also non-redundantly combines with negated predicates (see (\ref{ex:wiemer25})). These properties clearly indicate loss of morphological and semantic transparency, but, contrary to Macedonian Bulgarian \textit{da ne}, the modified clause needs to be introduced by another \textit{da} (see (\ref{ex:wiemer2}), (\ref{ex:wiemer17}), (\ref{ex:wiemer20}), (\ref{ex:wiemer25})).\footnote{There are exceptions to this, for instance if \textit{da ne bi} modifies a PP to express surprise, e.g.\ \textit{Pitam kelnera: ``A zeleto s kakvo go režete? –} \textit{\textbf{da ne bi} s~nožovka?''} `I ask the waiter: ``How do you cut the cabbage? - \textbf{da} \textbf{ne} \textbf{bi} \textbf{(${\approx}$ really)} with a hacksaw?''' \citep[148]{Ivanova2014}. Such mirative uses of \textit{da ne bi} are indicative of its lexicalization as an epistemic modifier, in analogy to Macedonian \textit{da ne} in (\ref{ex:wiemer34})--(\ref{ex:wiemer35}).} This subsequent \textit{da} can hardly be considered part of the complex marker itself, but must be analyzed as the usual verbal proclitic (sometimes within a clitic complex, as in (\ref{ex:wiemer2}), (\ref{ex:wiemer19}), and (\ref{ex:wiemer25})). In any case, for \textit{da ne bi} the irrealis requirement still applies, similarly to North Slavic units containing \textit{by}. This behaviour applies in main clauses as well, regardless of whether it marks apprehension (as in (\ref{ex:wiemer39})) or functions as an epistemic modifier (e.g., in rhetoric questions):\footnote{For details and slight differences between Macedonian \textit{da ne} and Bulgarian \textit{da ne bi} cf.\ \citet[147--150]{Ivanova2014} and \citet[70]{MitkovskaEtAl2017}.}


\ea 
\label{ex:wiemer39} 
\langinfo{Bulgarian}{}{\citealt[73]{MitkovskaEtAl2017}} \\
\gll Lele, da ne bi da Vi nastăp-ix po mazol-a? \\
\textsc{excl} \textsc{irr} \textsc{neg} \textsc{sbjv} \textsc{irr} \textsc{2pl.dat} step\textunderscore on\textsc{[pfv]-aor.1sg} over bunion\textsc{[m]-do} \\
\glt `Oh dear, have I maybe stepped on your bunion?' \\
\z



 Turning now again to North Slavic, the irrealis morpheme \textit{by} has become an integral part of clause-initial particles which cannot any longer be segmented into morphemes (see Footnote 13). If these connectives mark apprehension, negation must follow, although it need not be adjacent (compare, e.g., Russian \textit{kak by $
\ldots$ ne}, generalized as [X+\textit{by} \textsc{$\ldots$ neg}]). In semantic terms, too, the negation is transparent for those markers which have equivalents without negation in optative function (e.g., Russian \textit{liš' by} (\textit{ne}), \textit{kak by} (\textit{ne})). With respect to volitional attitude, optative and apprehensional meanings behave like antipodes: the judging subject (speaker or contextually determined) does not want E to occur (see \tabref{tab:componentsslavic} and \sectref{sec:conprem}; also \citealt[343]{Dobrushina2016}). This holds for independent and main clauses likewise. However, in embedded clauses the negation semantically harmonizes with the `fear'-component of the matrix predicate; we are thus dealing with ``apprehensional concord" (in analogy to modal concord); see (\ref{ex:wiemer40a}) below. On the one hand, concord phenomena support syntactic tightness between the relevant parts of the utterance (here a complex sentence); on the other hand, they imply that previously these parts were more autonomous from each other. As long as the assumed dependent part (here: the clause introduced by [X+\textit{by} \textsc{$\ldots$ neg}]) can occur on its own, concord simply testifies to usual semantic mechanisms of collocatability. (In \sectref{subsec:insubordination} we will see that this mechanism can be made responsible for embedding.)


Moreover, if in embedded apprehensional clauses the default complementizer (Russian \textit{čto}, Polish \textit{że}, Czech/Slovak \textit{že}) is used to denote an undesirable situation, negation is not required. If negation is employed, its use is compositional, i.e.\ it marks the judging subject's fear that the respective situation might not occur. Thus, (\ref{ex:wiemer40c}) is synonymous with (\ref{ex:wiemer40a}) only without negation (indicated by~\textsuperscript{\#}).



\ea 
\label{ex:wiemer40} 
\textnormal{Polish} \\
\ea \label{ex:wiemer40a} 
\gll Rob-imy razem interes-y \textnormal{($\ldots$)}
 i boj-ę się, \textbf{żeby} nie wpad-ł-\emptyset {} w zł-e towarzystw-o. \\
do\textsc{[ipfv]-prs.1pl} together deal\textsc{[m]-acc.pl} {} and fear\textsc{[ipfv]-prs.1sg} \textsc{refl} \textbf{\textsc{comp.irr}} \textsc{neg} fall\textsc{[pfv]-pst-3sg.m} in bad-\textsc{acc} company-\textsc{acc} \\
\glt `We are partners in business ($\ldots$) and I'm afraid that he might fall in with the wrong crowd.' {[PNC; M. Bielecki: \textit{Osiedle prominentów}. 1997}]

\ex \label{ex:wiemer40b}
\gll \textnormal{*($\ldots$)} \textbf{że} nie wpad-ł-\textbf{by}-\emptyset {} w zł-e towarzystw-o. \\
{} \textbf{\textsc{comp}} \textsc{neg} fall\textsc{[pfv]-pst-\textbf{sbjv}-3sg.m} in bad-\textsc{acc} company-\textsc{acc} \\
\ex \label{ex:wiemer40c}
\gll \textnormal{($\ldots$)} \textbf{że} (\textsuperscript{\#}nie) wpadni-e w zł-e towarzystw-o. \\
{} \textbf{\textsc{comp}} \textsc{neg} fall\textsc{[pfv]-fut.3sg} in bad-\textsc{acc} company-\textsc{acc} \\
\glt `($\ldots$) I'm afraid that he will (\textsuperscript{\#}not) fall in with the wrong crowd.'
\z
\z


Concomitantly, loss of morphological transparency in the univerbation of \textit{by} with other clause-initial particles has led to a situation in which complex connectives (X+\textit{by $\ldots$} \textsc{neg}) employed as negative-optative markers can be distinguished from free combinations of morphologically simple (and semantically neutral) connectives with the subjunctive marker (X + ($\ldots$ +) \textit{by}).\footnote{Cf.\ \citeauthor{Wiemer2015} (\citeyear{Wiemer2015}, \citeyear[\S3.1]{Wiemer2023}), also \citeauthor{Dobrushina2012} (\citeyear{Dobrushina2012}, \citeyear[322--327]{Dobrushina2016}) for Russian.} After `fear' verbs such combinations are ungrammatical (see (\ref{ex:wiemer40b})). However, at least in Polish we do find examples like (\ref{ex:wiemer41}) in which the standard complementizer (\textit{że}) introduces an apprehensional clause with an auxiliary complex and in which the subjunctive is marked on the auxiliary \textit{móc} `can'. Note, however, that an expletive negation is lacking and time-reference can only be posterior:


\ea 
\label{ex:wiemer41} 
\textnormal{Polish} \\
\gll Bartek obawia się, \textbf{że} mog-l-i-\textbf{by}-śmy tam wnieś-ć ze sobą nie najmądrzejsz-e pomysł-y z nasz-ego świat-a. \\
\textsc{pn.nom} fear\textsc{[ipfv].prs.3sg} \textsc{refl} \textbf{\textsc{comp}} can-\textsc{pst-pl-\textbf{sbjv}-1pl} there bring\textsc{[pfv]-inf} with \textsc{refl.ins} \textsc{neg} wisest-\textsc{acc.pl} idea-\textsc{acc.pl} from our-\textsc{gen.sg.m} world\textsc{[m]-gen.sg} \\
\glt `Bartek is afraid \textbf{that} there we \textbf{might} contribute not the best ideas of this planet.' {[PNC; D. Terakowska: \textit{Władca Lewawu}. 1989]}
\z


\noindent This pattern (standard complementizer + subjunctive of POSS auxiliary) resembles one of the options encountered in Serbian (see (\ref{ex:wiemer30c})).










\section{Negative-directive strategy (\textsc{NegDir})}
\label{sec:negdir}

The negative-directive strategy is realized by preventives.\footnote{On the history of the term `preventive' and its (quasi-)synonyms cf.\ \citet[29--30]{Ivić1958}; \citeauthor{Birjulin1992} (\citeyear[2]{Birjulin1992}, \citeyear[99--100]{Birjulin1994}).} Due to their directive illocutionary force, preventives cannot be embedded. This is a major difference from the negative-optative strategy \citep[41]{Dobrushina2006wiemer}: directive-preventives are used not just to express apprehension, but to alert the addressee to avoid the respective undesirable situation, i.e.\ the speaker does not want the addressee's actions (or lack thereof) to inadvertently cause such a situation (see \sectref{sec:conprem}). The inadvertent component can be emphasized by appropriate adverb(ial)s (see (\ref{ex:wiemer6}), (\ref{ex:wiemer42}), (\ref{ex:wiemer45b})):



\ea \textnormal{Polish} \\
\label{ex:wiemer42} 
\gll Nie zapomn-ij {\op}przypadkiem{\cp} okular-ów! \\
\textsc{neg} forget\textsc{[pfv]-imp.sg} accidentally glasses-\textsc{gen} \\
\glt `Don't forget (accidentally) your glasses!'
\z


\noindent Jointly with this, preventives are inherently future-oriented, and this further differentiates the \textsc{NegDir} strategy from \textsc{NegOpt} strategies.

Preventive utterances can include attention-catching devices. Many of them are imperatives, almost exclusively of imperfective verbs, e.g.\ Russian \textit{smotri} `look' (\citealt[153--154]{Volodin1986}, \citealt[215--216]{Xrakovskij1990}), Polish \textit{uważaj} `be careful', Czech \textit{dávej pozor} `pay attention', Slovene \textit{pazi} `be careful', Bulgarian \textit{pazi} \textit{se} `be careful', \textit{vnimavaj} `be attentive'. Only Russian \textit{ne} \textit{vzdumaj} ${\approx}$ `don't happen to think' derives from a perfective verb;\footnote{Other attention catchers attested in preventives are nouns like Polish \textit{uwaga}, Czech \textit{pozor} `attention' or particles like Czech \textit{ať}, which has various directive and purpose-related uses.} incidentally, it represents itself a (lexicalized) preventive imperative. These forms distinguish singular vs.\ plural in accordance with the number of addressees; they can therefore still be classified as imperatives and need not be treated as particles.




\subsection{Negated imperatives and the role of aspect}
\label{subsec:negimp}

Most accounts of directive-preventive strategies in Slavic are based on Russian, and it has been emphasized that the preventive function of the negated imperative is bound to perfective aspect. In fact, this function usually seems the only sensible interpretation of a negated perfective verb employed in a directive speech act. The central function of perfective aspect is to put a limit on the denoted eventuality, thus negation operates on a boundary, not on the action itself. Apart from preventives, this can be illustrated with sentences in which negation has narrow scope over some quantifiable component of the situation (\citealt[104]{Holvoet1989}; \citealt[38]{Luczkow1997}; \citealt[18]{Laskowski1998} for Polish; \citealt[122]{Kučera1984} for Czech), e.g.\ (\ref{ex:wiemer43}) and (\ref{ex:wiemer44}).

\newpage
\ea \textnormal{Polish} \\
\label{ex:wiemer43} 
\gll Nie zastaw \textbf{cał-ego} pokoj-u swo-imi grat-ami! \\
\textsc{neg} obstruct\textsc{[pfv].imp.sg} entire-\textsc{gen} room-\textsc{gen} \textsc{refl-poss-ins.pl} lumber-\textsc{ins.pl} \\
\glt `Don't obstruct the \textbf{whole} room with your lumber!'
\z


\ea \textnormal{Polish} \\
\label{ex:wiemer44}
\gll Nie wróć \textbf{zbyt} \textbf{późno!} \\
\textsc{neg} return\textsc{[pfv].imp.sg} too late \\
\glt `Don't return \textbf{too late!}' (i.e.\ `Return, but not too late!')
\z


 In preventives, negation acquires narrow scope since the purpose of the clause is only to indicate that a particular state-changing event (= boundary) should not be brought about, regardless of what the addressee might do to avoid it (see (\ref{ex:wiemer45b})). Thus the action itself is beyond its scope, or presupposed. By contrast, in prohibitives negation has wide scope, since the speaker urges the addressee not to undertake any activity connected to a state-changing event (see (\ref{ex:wiemer45a})). This contrast becomes particularly evident under minimal-pair conditions; compare the imperfective (\ref{ex:wiemer45a}) with the perfective (\ref{ex:wiemer45b}) negated imperative, both referring to a single time-located situation:



\ea \textnormal{Russian}
\label{ex:wiemer45} 
\ea \label{ex:wiemer45a}
\gll Ne otkryva-j okn-o! \\
\textsc{neg} open\textsc{[ipfv]-imp.sg} window-\textsc{acc} \\
\glt `Don't open the window!' \\
$\rightarrow$ prohibitive (IPFV) ${\cong}$ `Don't undertake any action connected to opening the window (e.g., turning the handle)!'
\ex \label{ex:wiemer45b}
\gll {\op}Smotri\textnormal{,)} ne otkro-j {\op}slučajno{\cp} okn-o! \\
look \textsc{neg} open\textsc{[pfv]-imp.sg} accidentally window-\textsc{acc} \\
\glt `(Look,) don't open the window (accidentally)!' \\
$\rightarrow$ preventive (PFV) ${\cong}$ `Whatever you do (e.g., with the handle), be careful not to open the window!'
\z
\z


\noindent The negated imperfective imperative yields a prohibitive meaning; the motivation for this is a volitional basis (`I don't want you to deal with situation S'). The preventive meaning, in contrast, is more closely associated with prediction (i.e.\ an epistemic function), which, in turn, often arises from inferences based on contingent circumstances (`If you don't take care, event E is going to happen'; see \tabref{tab:templatetwo} in \sectref{sec:conprem}). This correlates with other contexts of negation, for which a tendency toward a volition-based split of modality (deontic ${\rightarrow}$ \textsc{ipfv} vs.\ non-deontic ${\rightarrow}$ \textsc{pfv}) has been observed, at least in Russian and Polish.\footnote{Cf.\ \citeauthor{Wiemer2001} (\citeyear{Wiemer2001}, \citeyear[403--405]{Wiemer2008}, \citeyear[281--285]{Wiemer2017}) for overviews, \citeauthor{Padučeva2008} (\citeyear{Padučeva2008}, \citeyear[211--218]{Padučeva2013}) for Russian, and \citet{Divjak2009} for a usage-based investigation on Russian.} The common denominator is that negated perfective aspect triggers implicatures focusing on boundaries, while negated imperfective aspect can defocus boundaries; it is thus compatible with other parts of clausal semantics, for which implicatures arise. Thus, the distinction between prohibitive (${\rightarrow}$ \textsc{ipfv}) and preventive (${\rightarrow}$ \textsc{pfv}) readings of the negated imperative lines up well with more general tendencies of aspect choice, and these appear to be most salient in the northeastern part of Slavic (= Russian). Inner-Slavic differences result from the different strength with which invited inferences have conventionalized and been associated with perfective and imperfective aspect in complementary distribution.


We may thus say that a preventive interpretation arises from the strengthening of implicatures resulting from a transparent combination of negation and perfective aspect in directive illocutions for time-located events. This interpretation is entirely compositional. Nonetheless, the conventionalization of invited inferences, based on the binary \textsc{pfv:ipfv} opposition, appears to be weaker in South Slavic, even if we account for constructions that have been ousting the inherited imperative (see \sectref{subsec:analyticnegimp}). Moreover, the strict aspect-based division between prohibitive and preventive is definitely a rather recent development even in Russian (see \sectref{subsec:preventives}). 


Simultaneously, we have to realize that a person uttering a preventive speech act must assume that the addressee has control~– not necessarily over the event denoted by the verb in the negated imperative, but over associated actions which can allow the addressee to avert the undesirable event (or its consequences); see \tabref{tab:templatetwo} in \sectref{sec:conprem}. Otherwise we could hardly explain why preventives are frequent with controllable actions (see (\ref{ex:wiemer45b}) and (\ref{ex:wiemer47a})), but excluded for certain verbs, e.g.\ of mental activity, for which control cannot be given a limit; compare Polish   *\textit{Nie zastanów}\textsc{\textsuperscript{pfv}} \textit{się nad tym!} \citep[103]{Holvoet1989}, to mean something weird like `Don't (make it happen that you) accidentally think it over'. Importantly, although preventive utterances often build on non-agentive verbs, agentive verbs are not excluded (see (\ref{ex:wiemer45b}) and Russian \textit{Ne slomaj}\textsc{\textsuperscript{pfv}} \textit{stul} `Don't break the chair', \textit{Ne razbej}\textsc{\textsuperscript{pfv}} \textit{steklo} `Don't break the glass'). Preventives ``denote uncontrollable events resulting from some relevant controllable actions, and it is these controllable actions that are in fact addressed by the prescriptor (e.g.\ the sentence \textit{Don't break the chair} may be interpreted as `Don't rock the chair' [so that it will not break; BW])'' (\citealt[38]{BirjulinXrakovskij2001}; also \citealt{Xrakovskij1990}, \citealt{Birjulin1992}).


Negated perfective imperatives happen to acquire other pragmatic functions as well, e.g.\ advice, request, or persuasion (e.g., Polish \textit{Nie obraź}\textsc{\textsuperscript{pfv}} \textit{się na mnie} `Don't get offended at me', Serbian\ \textit{Ne shvati}\textsc{\textsuperscript{pfv}} \textit{krivo!} `Don't misunderstand me!') or, to the contrary, categorical demands and pleas. Such functions are ``fuzzier'' than the preventive; they partially depend on the lexically determined actionality of the verb, partly they can be calculated from specific social and conversational constellations between the interlocutors (who is in control of what and over whom?).\footnote{Cf.\ the probably still most subtle and well-considered analyses to date in \citet[101--106]{Holvoet1989} and \citeauthor{Birjulin1992} (\citeyear{Birjulin1992}, \citeyear[\S4]{Birjulin1994}). Cf.\ also \citet[36--37]{Ivić1958}, \citet[22--26]{Laskowski1998}.} Many such instances are archaic and have become petrified phrases (e.g., Russian/Bulgarian \textit{Ne daj Bože} `God forbid', lit.\ `Don't give, God'). Here belong some of the aforementioned attention catchers (e.g.\ Russian \textit{ne vzdumajte} `don't (dare) think'), and we find them also in varieties of South Slavic which have otherwise abandoned perfective verb forms in negated imperatives; e.g.\ Serbian \textit{Ne uzmi}\textsc{\textsuperscript{pfv}} \textit{za zlo} `Don't take it for harm', against *\textit{Ne uzmi}\textsc{\textsuperscript{pfv}} \textit{ovu knjigu} `Don't take (inadvertently) this book' or *\textit{Ne nazebi}\textsuperscript{\textsc{pfv}} `Don't get a cold' (\citealt[27--28]{Ivić1958}), which would be truly preventive, but have been ousted by periphrastic constructions (see \sectref{subsec:analyticnegimp}).

Moreover, negated perfective imperatives occur throughout Slavic as a means of threat or intimidation, although in South Slavic (except Slovene) this function is residual (see \sectref{subsec:analyticnegimp}). Compare, for instance, Czech \textit{Jenom mi to} \textbf{\textit{nekup}}!\textsuperscript{\textsc{pfv}} `\textbf{Don't (you) dare buy} it for me' (\citealt[28--29]{Ivić1958}; \citealt[60--61]{Kopečný1962}), Macedonian \textbf{\textit{Nedojdi}\textsc{\textsuperscript{pfv}}} \textit{pa} \textit{\textsc{\'k}e} \textit{vidiš što te čeka!} `\textbf{Don't come}, or you'll see what you get.' (lit.\ `$\ldots$ what waits') (\citealt[416]{Koneski1976}, \citealt[346--347]{Tomić2012}, \citealt{Mitkovska2015}).  Obviously, a warning is the verbalization of a threat, and threats are conceptually very close to preventives; the crucial difference is that threats are uttered in the interest of the speaker, while preventive speech acts are issued for the sake of the addressee (see \sectref{sec:conprem}).


After all, we can say that, if negated perfective imperatives have not been ousted as such in the respective Slavic language, their preventive function predominates across the verb lexicon; other, minor functions are related to preventives and can be explained via pragmatic mechanisms. As for those Slavic languages in which negated perfective imperatives are no longer employed (basically, South Slavic except Slovene), it is an open question whether equivalent periphrastic constructions have taken over all other functions of negated perfective imperatives as well (see \sectref{subsec:analyticnegimp}).







\subsection{Areal cline of functions ruled by aspect?}
\label{subsec:areal}


It has been argued that in comparison to Russian (and Polish), other Slavic languages exhibit a less clear-cut prohibitive-preventive distinction ruled by aspect.\footnote{Cf.\ \citet[96]{Nilsson2013} in general and \citet[156]{Ivanova2014}, \citet{MitkovskaEtAl2017} for \textit{da}-equivalents of negated imperatives in Bulgarian and Macedonian.} Such qualifications require differentiation. Let us look at the western periphery of Slavic, before we turn to the central and eastern part of South Slavic (see \sectref{subsec:analyticnegimp}).


First, we do find negated perfective imperatives in preventive function in Slovene, as in (\ref{ex:wiemer46}).


\ea \textnormal{Slovene} \\
\label{ex:wiemer46} 
\gll Ne prehlad-i se! \\
\textsc{neg} get\textunderscore cold\textsc{[pfv]-imp.sg} \textsc{refl} \\
\glt `Don't get a cold!'
\z



\noindent The question is how consistent this use (and its distinction from prohibitives) is in Slovene. Empirical research is lacking, however an educated guess is that consistency is much lower than in Russian (M. Uhlik, p.c.).

Second, Czech employs negated perfective imperatives in preventive meanings as well, and we do find an opposition to prohibitive meaning with imperfective verbs (\citealt[322]{Trávníček1923}, \citealt[79--81]{Dokulil1948}, \citealt[60--61]{Kopečný1962}, \citealt{Kučera1984}; also \citealt[366--367]{Meyer2010} for Slovak), as in (\ref{ex:wiemer47}).


\ea \textnormal{Czech}
\label{ex:wiemer47}
\ea \label{ex:wiemer47a}
\gll Ne-škrábn-i se! \\
\textsc{neg}-scratch\textsc{[pfv]-imp.sg} \textsc{refl} \\
\glt `Do not scratch yourself accidentally!'
\ex \label{ex:wiemer47b}
\gll Ne-škrab se! \\
\textsc{neg}-scratch\textsc{[ipfv].imp.sg} \textsc{refl} \\
\glt `Do not scratch yourself (any more)'; `Stop scratching!'
\z
\z


 The distinction can also be observed with verbs whose aspect correlates with different ontological types of objects. This, in turn, conditions different actionality interpretations of the entire clause; compare (\ref{ex:wiemer48a}) and (\ref{ex:wiemer48b}).


\newpage´
\ea \textnormal{Czech}
\label{ex:wiemer48} 
\ea \label{ex:wiemer48a}
\gll Ne-ztrať ten klíč! \\
\textsc{neg}-lose\textsc{[pfv].imp.sg} this key\textsc{.acc.sg} \\
\glt `Do not lose the key!'
\ex \label{ex:wiemer48b}
\gll Ne-ztráce-j naděj-i! \\
\textsc{neg}-lose\textsc{[ipfv]-imp.sg} hope-\textsc{acc.sg} \\
\glt `Do not lose hope!'
\z
\z


 However, the distribution of perfective and imperfective verbs in the negated imperative seems to follow from a more general tendency due to which in Czech aspect choice is more consistently governed by different actionality types than in Russian. In the latter, imperfective verbs have expanded into contexts for which languages in the western half of the Slavic-speaking world rather choose perfective aspect. For instance, ``Czech utilizes perfectives in negative constructions more widely than Russian does, and admits the imperfective in negative event descriptions much more rarely'' \citep[118]{Kučera1984}. This observation fits well with the west-east cline of Slavic aspect commonly known since \citet{Dickey2000}. Nonetheless, the preventive-prohibitive pattern which we discover for Czech, Russian and Polish must be based on the same considerations which were pointed out in \sectref{subsec:negimp}. \citet{Kučera1984} based his considerations on the factor of control: negated imperative and [+control] favours imperfective, negated imperative and [−control] favours perfective aspect. But these considerations relate control only to the lexical meaning of the verb used in the imperative; they do not account for the meaning of the whole construction, which implies that the speaker wants the addressee to undertake measures (see \tabref{tab:templatetwo} in \sectref{sec:conprem}). If we take this into account, Kučera's analysis confirms the analysis proposed in \sectref{subsec:negimp}, and we may conclude that the actionality-based preferences for perfective or imperfective verbs~– which are generally stronger in Czech than in Russian~– provide a favorable basis for the aspect-driven contrast between prohibitive and preventive interpretations in negated directive speech acts.


\subsection{Analytic negated imperatives in South Slavic}
\label{subsec:analyticnegimp}

In Balkan Slavic, negated imperatives of perfective verbs have receded, but for independent reasons. These cannot be discussed here, but it is intriguing to observe that the retreat of perfective verbs from negated imperatives has been accompanied by an increase of periphrastic constructions based not only on \textit{da}, but also on semi-petrified auxiliary forms with incorporated negation: Bulgarian \textit{nedej/nedejte} (\textsc{sg/pl}) (+ \textit{da}) < \textit{ne děj(te)} from Common Slavic \textit{dějati} `do', Serbian \textit{nemoj/nemojmo/nemojte} (\textsc{2sg/1pl/2pl}), Macedonian \textit{nemoj/nemojte} (\textsc{2sg/2pl}) <\textit{ne mogi} from \textit{moći} `can' (\citealt{Ivić1958}, \citealt{Sedláček1997}). Note that these auxiliaries still distinguish number, as do attention catchers derived from verbs (see \sectref{sec:negdir}).


Bulgarian \textit{nedej} combines with imperfective verbs in a form which goes back to the truncated infinitive, but is now homonymous with the \textsc{2/3.sg} aorist, e.g.\ \textit{Nedej(te) pisa!} `Don't write!' (\citealt{Ivić1958}, \citealt{Sedláček1997}, \citealt[412]{Lindstedt2010}). Important is the restriction to imperfective verbs, so that this construction blocks a prohibitive-preventive opposition; prohibitive meanings are the default. By contrast, Serbian/Macedonian\ \textit{nemoj} combines with either perfective or imperfective aspect; Macedonian \textit{nemoj} combines only with \textit{da}+finite verb (since even the truncated infinitive is lost), while its Serbian cognate optionally allows for the infinitive. However, aspect choice does not yield a reliable opposition of prohibitive (${\rightarrow}$ \textsc{ipfv}) vs.\ preventive (${\rightarrow}$ \textsc{pfv}) function. Therefore, the four possible realizations in (\ref{ex:wiemer49}) are practically synonymous; in combination with imperfective verbs, in any case, they trigger a prohibitive reading, while with a perfective verb no straightforward preventive reading arises (for Macedonian equivalents cf.\ \citealt[345]{Tomić2012}):\footnote{See, however, \citet[400]{Szucsich2010} on connections between an apprehensional (`threat') reading and the choice of \textit{da}+V\textsubscript{fin} or infinitive after \textit{nemoj}. He also notes that \textit{nemoj(-mo/-te)} need not always  agree with the finite verb after \textit{da}.}


\ea \textnormal{Serbian}
\label{ex:wiemer49} 
\ea \label{ex:wiemer49a} \textnormal{\textit{nemoj} + pfv. infinitive} \\
\gll Nemoj sada otvori-ti prozor! \\
\textsc{neg.aux.2sg} now open\textsc{[pfv]-inf} window\textsc{.acc.sg} \\
\ex \label{ex:wiemer49b} \textnormal{\textit{nemoj} + ipfv. infinitive} \\
\gll Nemoj sada otvara-ti prozor! \\
\textsc{neg.aux.2sg} now open\textsc{[ipfv]-inf} window\textsc{.acc.sg} \\
\ex \label{ex:wiemer49c} \textnormal{\textsc{neg} + ipfv. infinitive} \\
\gll Ne otvara-j sada prozor! \\
\textsc{neg} open\textsc{[ipfv]-imp.sg} now window\textsc{.acc.sg} \\
\ex \label{ex:wiemer49d} \textnormal{\textit{nemoj} + \textit{da} + pfv./ipfv. present} \\
\gll Nemoj sada da otvori-š\textnormal{/}otvara-š prozor! \\
\textsc{neg.aux.2sg} now \textsc{irr} open\textsc{[pfv]-prs.2sg}/open\textsc{[ipfv]-prs.2sg} window\textsc{.acc.sg} \\
\glt `Don't open the window now!'
\z
\z


 Perfective verbs in negated imperatives have more and more been driven out by periphrastic constructions, so that the original imperative (see (\ref{ex:wiemer49c})) no longer has a perfective equivalent. A preventive reading in periphrastic constructions becomes more likely with low-agentivity verbs (J. Grković-Major, p.c.):


\ea \textnormal{Serbian} \\
\label{ex:wiemer50} 
\gll Nemoj da se razboli-š! \\
\textsc{neg.aux.2sg} \textsc{irr} \textsc{refl} get\textunderscore ill\textsc{[pfv]-prs.2sg} \\
\glt `Don't get ill!'
\z


 Apart from ousting negated perfective imperatives, these periphrastic constructions correspond to techniques which have replaced the infinitive. Infinitive loss is commonly known as a Balkan feature, and within South Slavic we observe a clear (south)east~-(north)west cline, with standard Serbian and Croatian occupying an intermediate position. The most common replacement technique is the irrealis marker \textit{da} which combines with the present tense or the perfect. This construction also replaces the negated perfective imperative, e.g.


\ea \textnormal{Bulgarian}\\
\label{ex:wiemer51} 
\gll {\op}Vnimava-j,{\cp} da ne zabravi-š knig-a-ta!\\
pay\textunderscore attention\textsc{[ipfv]-imp.sg} \textsc{irr} \textsc{neg} forget\textsc{[pfv]-prs.2sg} book\textsc{[f]-sg-def.sg.f}\\
\glt `(Attention) Don't forget the book!'
\z

\ea 
\label{ex:wiemer52} \textnormal{(= \ref{ex:wiemer37})} \\
\langinfo{Macedonian}{}{\citealt[30]{Mitkovska2015}} \\
 \gll {\op}Pobrza-jte!{\cp} Da ne go ispušti-te avion-ot! \\
hurry\textunderscore up\textsc{[pfv]-imp.pl} \textsc{irr} \textsc{neg} \textsc{3m.acc} miss\textsc{[pfv]-prs.2pl} airplane\textsc{.sg.m-def.sg.m} \\
\glt `(Hurry up!) Don't miss the plane.' \\
or: `$\ldots$ for you in order not to miss the plane.' (connection with purpose clauses) \\
\z

\noindent Superficially, this construction may be difficult to distinguish from the negative-optative strategy (see \sectref{subsec:negoptcomponents}). As in Serbian, preventive meaning is claimed to obtain only if the denoted event is beyond the addressee's control (\citealt[41]{Cakarova2009}). Compare Serbian in (\ref{ex:wiemer50}) with Bulgarian in (\ref{ex:wiemer53}):

\ea \textnormal{Bulgarian} \\
\label{ex:wiemer53} 
\gll Da ne si razvali-š stomax-a! \\
\textsc{irr} \textsc{neg} \textsc{refl.dat} destroy\textsc{[pfv]-prs.2sg} stomach-\textsc{do} \\
\glt `Don't (accidentally) spoil your stomach.'
\z


\noindent The claim concerning the addressee's control requires the same corrections as for North Slavic (see \sectref{subsec:negimp}). After all, we come to the preliminary conclusion that not all South Slavic ``analytic imperatives'' allow for perfective verbs, but if they do, aspect choice does not yield a reliable distinction of prohibitive vs.\ preventive function to the same degree as in Slavic languages further to the north and east. Moreover, a preventive reading of a perfective negated imperative appears to be more likely if it does not call precautional actions by their name.




\section{Diachronic background}
\label{sec:diachronic}

The diachronic background of apprehensional strategies is largely a \textit{terra incognita}. Therefore, the following considerations are tentative, restricted to basic issues, and should rather set the stage for future research.


\subsection{Negative-optative strategy resulting from insubordination?}
\label{subsec:insubordination}

Negative-optative and negative-purpose clauses can be reinterpreted as complement clauses. However, even in cases when such clauses immediately follow predicates with an apprehensional semantics, their complement status often remains debatable and they show signs of low syntactic integration (see \sectref{subsubsec:negoptnorthslavic}). From a synchronic perspective, clauses with apprehensional markers in initial position may sometimes look like truncated sentences, i.e.\ with an elliptic main clause (containing the matrix predicate). This would amount to insubordination, i.e.\ the ``conventionalized main clause use of what, on prima facie grounds, appear to be formally subordinate clauses'' \citep[367]{evans2007insubordination}. Evans himself is aware that such a \textit{prima facie}-qualification can turn out to be inadequate after an assessment of the diachronic background. That is, what superficially resembles the remainder of a complex sentence (with elided main clause) can actually turn upside down the chronology, which rather testifies to a development from juxtaposition to clausal complementation.


Consequently, the question to be asked is: did Russian \textit{kak by $\ldots$  ne}, Polish \textit{żeby $\ldots$  nie}, Bulgarian \textit{da ne bi}, etc., first occur as complementizers (i.e.\ as flags indicating that the following clause functions as an argument of a predicative expression in the matrix clause), or did they simply occur as clause-initial units able to mark apprehension? I want to argue that, in general, diachronic primacy should be assigned to usage in independent clauses, i.e.\ that juxtaposition preceded tighter clause combining. While this eventually might not apply to constructions deriving from negative purpose clauses (see \sectref{subsubsec:negpurpose}), for all other apprehensional clause types discussed here we find strong conceptual and empirical arguments against insubordination.


\subsubsection{Negative-optative in North Slavic}
\label{subsubsec:negoptnorthslavic}


A plausible diachronic explanation for the rise of embedded negative-optative clauses looks as follows. Independent clauses introduced with negative-optative markers frequently occurred after `fear' verbs (or related predicative units). The latter were suitable as complement-taking predicates (CTPs) of clauses headed by negative-optative markers, so that the latter became reinterpreted as clausal complements of predicates because they ``harmonized'' with them in terms of apprehensional semantics. This process can be called ``syntactic tightening by apprehensional harmony" (cf.\ \citealt[54]{Nilsson2012} and \citealt[345--346]{Dobrushina2016} for a similar reasoning). This is shown in \figref{tab:syntactictightening} for Russian \textit{kak by $\ldots$  ne}, illustrated in (\ref{ex:wiemer54}).





\ea 
\label{ex:wiemer54} 
\langinfo{Russian}{}{RNC; S. Taranov 1999; cited in \citealt[59]{Nilsson2012}}\\
\gll Ja \textit{boja-l-\emptyset-sja}, \textbf{kak} \textbf{by} ty \textbf{ne} zabole-l-a posle svo-ix skitani-j. \\
\textsc{1sg.nom} fear\textsc{[ipfv]-pst-sg.m-refl} how \textsc{irr} \textsc{2sg.nom} \textsc{neg} get\textunderscore ill\textsc{[pfv]-pst-sg.f} after \textsc{refl.poss-gen.pl} wandering-\textsc{gen.pl} \\
\glt `I \textit{was afraid} \textbf{that} you became/you could/might/would become ill after all your wanderings.' (Zorikhina Nilsson's translation)
\z

\begin{figure}[t]

\caption{Syntactic tightening of ``apprehensionally harmonious'' juxtaposed clauses}
\label{tab:syntactictightening}
\begin{tabular}{lllll}
\textit{}  & \textit{\textbf{p}} & \textit{\textbf{q}} &  & \\

\midrule

\textbf{{[}Stage A{]}} & X fears : & \textit{kak by  $\ldots$} & \textit{ne} & + V$_{fin}$/V$_{infinitive}$ \\
  &  &   {clause-initial particle} & \textsc{neg} & \\
\\
& \multicolumn{4}{l}{$\rightarrow$ juxtaposition (\textit{q} is related to \textit{p} at discourse level)} \\
\\

\textbf{{[}Stage B{]}} & X fears & \textit{kak by  $\ldots$} & \textit{ne} & + V$_{fin}$/V$_{infinitive}${]} \\
                 & CTP     & \multicolumn{2}{l}{complementizer (semantically} & \\
                 &    &  {``harmonic'' with CTP)} & \\
              \\
& \multicolumn{4}{l}{$\rightarrow$ complementation (\textit{q} becomes embedded in \textit{p})}
\end{tabular}
\end{figure}




 There are quite a few cases in the history of various, largely unrelated languages in which apprehensional complements are based on ``apprehensional harmony" (cf.\ \citealt{dobrushina2021} for a critical survey). As for Russian \textit{kak by $\ldots$ ne}, this complex connective appeared in complement clauses only in the 19\textsuperscript{th} century (\citealt[140--142]{dobrushina2021}). In independent clauses the negation is transparent inasmuch as it marks the negative equivalent of an optative clause, but in clausal complementation negation is simply in concord with the meaning of the apprehensional matrix predicate (see \sectref{sec:negopt}). I see no other way of explaining this relation than as a result of a development from juxtaposition to subordination.


Further support for the process assumed in \figref{tab:syntactictightening} comes from the following considerations. First of all, not all clause-initial optative and apprehensional markers occur in embedded clauses (e.g., Russian \textit{liš' by ($\ldots$ ne)}, Polish \textit{oby~\mbox{($\ldots$ nie)}}), but all complementizers of embedded apprehensional (or, respectively, optative) clauses occur in independent clauses. This directed implicational relationship speaks in favour of the assumption that the diachronic source of complementizers resides in clause-initial particles, not vice versa. The question, then, is which of the relevant clause-initial markers eventually gets involved in clausal complementation, as illustrated in (\ref{ex:wiemer55}) for example (\ref{ex:wiemer54}).




\newcommand{\tikzmarkmine}[1]{\tikz[remember picture] \node (#1) {};}
\eabox{
\label{ex:wiemer55}
%\begin{table}
\setlength{\tabcolsep}{1pt}
\begin{tabular}{lll}
\textit{\textbf{p}} & \multicolumn{2}{c}{\textit{\textbf{q}}}\\
 & \tikzmarkmine{A}\ding{51}\textit{kak by}\\
 \textnormal{X is afraid}& {~~* }\textit{tol'ko by}  & \textit{\hspace{3pt} $\ldots$ ty ne zabolela posle svoix skitanij}\\
 & {~~* }\textit{liš' by}\\
 & \tikzmarkmine{C}\textit{$\cdots$}\\
 &  \textnormal{`that/lest}       &  \textnormal{\hspace{3pt} $\ldots$ you become ill after all your wanderings.'}\\
    & \multicolumn{2}{l}{\begin{tabular}[l]{@{}l@{}} \textnormal{(Compare as independent clause: `If only you don't get ill $\ldots$')}\end{tabular}}
 \begin{tikzpicture}[remember picture, overlay]     
 \draw[middle color=gray, fill opacity=0.3, rounded corners]
        ($(A.north west)+(0.1,0.2)$) -- 
        ($(A.north east)+(1.7,0.2)$) --  
        ($(C.south east)+(1.7,-0.1)$) --  
        ($(C.south west)+(0.1,-0.1)$) -- 
        cycle;
     \end{tikzpicture}
\end{tabular}
%\end{table}
}

Finally, syntactic subordination between \textit{p} and \textit{q} remains weak. First, \textit{q} can be separated from \textit{p} by a pause (and a falling intonation contour at the end of \textit{p}). Second, \textit{p} and \textit{q} cannot be permuted. Moreover, I am unaware of any examples in which \textit{q} appears inserted into \textit{p}. Third, island extraction of a constituent from \textit{q} by a WH-word sounds very unnatural (and possibly does not exist); compare Russian \textsuperscript{??}\textit{Čego ty boiš'sja, kak by ne slučilos'?}, to mean `What are you afraid of that might happen?' For most of these properties cf.\ \citeauthor{Dobrushina2016} (\citeyear[344--345]{Dobrushina2016}), who also points out that they are characteristic of positive optative clauses (with \textit{kak by}) as well.



\subsubsection{Negative purpose in North Slavic}
\label{subsubsec:negpurpose}


\noindent The insubordination issue differs for connectives that introduce clauses whose structure demonstrates a close relationship to negative purpose. The diachronic relation between negative purpose clauses and complements to `fear' verbs in Slavic leaves many open questions, which cannot be solved here, but this relation is evident for clauses embedded by clause-initial particles which are composed of a standard complementizer (Russian \textit{čto}, Polish \textit{że}) or a coordinative conjunction (Polish/Czech/Slovak \textit{a}), the irrealis morpheme \textit{by} (Russian \textit{čtoby}, Pol. \textit{żeby}, Czech/Polish \textit{aby}) and clausal negation.

In the contemporary languages the aforementioned connectives are widely used as (main or almost exclusive) adverbial subordinators of (positive and negative) purpose clauses. Only Russian \textit{čtoby $\ldots$ ne} has become obsolete in apprehensional function, while the other ones are still well{}-attested in contemporary usage. Example (\ref{ex:wiemer56}) illustrates such an obsolete use:

\ea 
\label{ex:wiemer56} 
\langinfo{Russian}{}{A.\ Chekhov; cited in \citealt[55]{Nilsson2012}} \\
\gll \ldots{} \textbf{iz} \textbf{strax-a}, čtoby o nem ne duma-l-i durno. \\
{} out.of fear-\textsc{gen} \textsc{comp.irr} about \textsc{3m.loc} \textsc{neg} think\textsc{[ipfv]-pst-pl} badly \\
\glt `(He never kept maidservants) \textbf{fearing/out of fear} that they could/might/would think ill of him.' \\
\z

 Contrary to optative clauses and clauses with an apprehensional semantics, purpose clauses can hardly occur as independent clauses,\footnote{Cf.\ \citet[271--272]{Dobrushina2016}. This corresponds to a usual typological pattern \mbox{(cf.\ \citealt[150--151]{Schmidtke-Bode2009}).}} so that one wonders whether independent apprehensional clauses with the aforementioned clause-initial connectives are not to be considered instances of insubordination. Embedded clauses with an apprehensional meaning (after suitable CTPs) are often indistinguishable from negative purpose clauses (see (\ref{ex:wiemer13}), (\ref{ex:wiemer22}), (\ref{ex:wiemer37})). From a typological perspective, ``the evidence for a close structural relationship between purpose and complement clauses is overwhelming'', and there are several functional motivations for an overlap of purpose and certain complementation relations (\citealt[161, 164]{Schmidtke-Bode2009}). Nonetheless, the relation of negative purpose (>~apprehensional) to purpose is less clear than it might appear at first sight.\footnote{First of all, negative purpose clauses disfavour contexts with (directed) movement, while positive purpose clauses occur predominantly in exactly this context. Moreover, negative purpose clauses prefer different-subject constructions, whereas positive purpose clauses prefer same-subject constructions. These differences are indicative of different experiential correlations (cf.\ \citealt[129--144]{Schmidtke-Bode2009} for typological facts and a critical survey); probably, these correlations correspond to different illocutionary purposes.}


This leaves us with the question whether the apprehensional usage is diachronically secondary, i.e.\ derived from negative purpose\footnote{In Ancient Greek \textit{mé:} was used in apprehensional complement clauses before negative purpose clauses with the same subordinator appeared on the scene (\citealt[142--143]{dobrushina2021}, referring to \citealt{Goodwin1897}). The diachronic relation between these two clause types can thus be reversed, which calls for a case-by-case investigation (prior to generalizations about the direction).} (see (\ref{ex:wiemerDiachronya})), or whether both functions developed in parallel, but independently from some other common source (see (\ref{ex:wiemerDiachronyb})).


\ea \label{ex:wiemerDiachrony}
\ea \label{ex:wiemerDiachronya}
\begin{tabular}{lllll}
{\ldots} \tikzmark{4} && \tikzmark{5} \textnormal{negative purpose} \tikzmark{6} && \tikzmark{7} \textnormal{apprehensional} \\
\tikz[remember picture] \draw[-latex, overlay] ([shift={(0.1em, 0.7ex)}]pic cs:4) -> ([shift={(-0.1em, 0.7ex)}]pic cs:5);
\tikz[remember picture] \draw[-latex, overlay] ([shift={(0.1em, 0.7ex)}]pic cs:6) -> ([shift={(-0.1em, 0.7ex)}]pic cs:7);
\end{tabular}
\ex \label{ex:wiemerDiachronyb}
\begin{tabular}{lll}
&& \tikzmark{1} \textnormal{negative purpose}\\
?? \tikzmark{2} \\
&& \tikzmark{3} \textnormal{apprehensional}
\tikz[remember picture] \draw[-latex, overlay] ([shift={(0.1em, 0.7ex)}]pic cs:2) -> ([shift={(-0.1em, 0.7ex)}]pic cs:1);
\tikz[remember picture] \draw[-latex, overlay] ([shift={(0.1em, 0.7ex)}]pic cs:2) -> ([shift={(-0.1em, 0.7ex)}]pic cs:3);
\end{tabular}
\z
\z

 As for Russian \textit{čtoby $\ldots$ ne}, \citet[142--143]{dobrushina2021} firmly states ``that independent wish-constructions were not the source of fear-complements with \textit{čtoby} (as they were with \textit{kak by} clauses); independent wish-clauses with \textit{čtoby} were quite rare, and all of the attested examples are instantiations of insubordination (see \citealt[337--346]{Dobrushina2016}).'' Evidently, the history of apprehensional complements with \textit{čtoby $\ldots$ ne} and \textit{kak by $\ldots$  ne} differed not only with respect to insubordination (yes for \textit{čtoby $\ldots$  ne}, no for \textit{kak by $\ldots$ ne}), but also in terms of their age and their further distributional behavior: \textit{čtoby $\ldots$  ne} in apprehensional complements appeared in Russian since the 17th century, it preceded \textit{kak by $\ldots$  ne} in this function by 200 years (\citealt[142--143]{dobrushina2021}). Then \textit{čtoby $\ldots$  ne} and \textit{kak by $\ldots$  ne} co-existed during the 19th century, also in independent clauses \citep[341]{Dobrushina2016}, but after the early 20th century \textit{čtoby $\ldots$  ne} started becoming obsolete as an apprehensional marker both in complement and in independent clauses (\citealt[58--60]{Nilsson2012}). It is now the practically exclusive adverbial subordinator for purpose clauses, while \textit{kak by $\ldots$  ne} is restricted to independent and complement apprehensional clauses. Russian has thus approached a stage of an almost complementary division between these two connectives, contrary to their equivalents in West Slavic, which are used in either function (which can often hardly be discriminated in particular discourse tokens). However, diachronic research is required to state with more certainty the direction of development (in terms of (\ref{ex:wiemerDiachrony})) and to give a qualified judgment on the insubordination issue. Among other things, we need to know when the relevant clause-initial connectives started coalescing with \textit{by} and how this process was related to their employment in purpose and wish clauses (and their negative equivalents).


\subsubsection{Balkan Slavic \textit{da ne}}
\label{subsubsec:balkanslavic}

Analogous reasoning applies, \textit{ceteris paribus}, to clause-initial units employed in apprehensional contexts in South Slavic languages. As for Balkan Slavic, we saw that the association with apprehension has remained much looser than in North Slavic (see \sectref{sec:negopt}). Macedonian/Bulgarian \textit{da ne} and Bulgarian \textit{da ne bi} reveal a high degree of context--dependence for their interpretation, with apprehension  usually as only one among many readings. Most probably, the contemporary situation here continues the initial vagueness of \textit{da}-clauses which was, as it were, inherited by \textit{da ne}-clauses even if this complex unit has become opaque to some degree (see \sectref{subsec:lossofcomp}). In particular, whether embedded apprehensional clauses derive from negative purpose or negative optative clauses, or rather arose directly under the scope of `fear' predicates, or whether different lines of development converged, remains an open question since no diachronic research has been conducted and reliable historical evidence is lacking. Cf.\ \citet[76--77]{MitkovskaEtAl2017}, who also can only surmise that Bulgarian \textit{da ne bi} could have ``originated from an optative source via the combination of the inherited subjunctive \textit{bi} and the Balkan subjunctive \textit{da}''.


\subsubsection{Summary on \textsc{NegOpt} strategies}
\label{subsubsec:summarynegopt}


All relevant clausal connectives are based on irrealis morphemes with concomitant restrictions of TAM marking (North Slavic \textit{by}, South Slavic \textit{da}), even if morphological segmentability has been lost partially or entirely.\footnote{The irrealis-feature of \textit{da} has been lost in the western part of South Slavic, but the rise of relevant clause-initial markers considerably predated this process (\citealt[295--303]{Wiemer2018}, with further references).} All these units lower (or suspend) assertiveness. Lowered assertiveness has been considered a typical feature of dependent clauses; however we also find it in all sorts of self-standing utterances with directive or optative illocutions, and such utterances need not be understood as descendants of more complex structures. It does not surprise, then, to find such self-standing clauses also for apprehension. We can safely assume that almost all etymologically complex units with \textit{da} and \textit{by}, respectively, started to head clauses with an apprehensional semantics prior to becoming reanalyzed as complementizers after suitable complement-taking predicates. The only exception is probably Russian \textit{čtoby $\ldots$  ne}, whose employment in apprehensional complements is derived from negative purpose clauses and employment in independent apprehensional clauses was rare (see \sectref{subsubsec:negpurpose}); for its West Slavic equivalents this issue needs to be investigated.

Thus, with this one exception, no ``main clause ellipsis'' need be assumed (\textit{pace} \citealt{Dobrushina2006wiemer} and \citealt{evans2007insubordination}); rather, the clause-initial employment of almost all apprehensional markers in independent clauses can be regarded as diachronically primary. This alternative analysis allows the appearance of embedded negative-optative clauses to be anchored to a unified basis which explains the rise of complex sentences from discourse pragmatics and, possibly, from the frequency of juxtaposition with suitable, ``semantically harmonious'' predicative expressions. In fact, this explanation for the appearance of embedded negative-optative clauses has (at least) two advantages over assuming insubordination.

First, it complies with the common-sense assumption that syntactic complexity (namely, tightening) increases over time. Of course, acquired complexity can again be reduced, e.g.\ by insubordination. However, if the question is about diachronic primacy (what comes first?), diachronic data often shows that alleged cases of insubordination turn out to be independent clause patterns which existed from their earliest attestations. This brings us to the next point.

Second, my explanation is in keeping with analyses of the rise of clausal complementation for other semantic classes of predicates (i.e.\ potential CTPs). Compare, for instance, Old Serbian \textit{da} \citep[198]{Grković-Major2004} and Latvian \textit{lai} (\citealt[237, 253, 260]{Holvoet2016}), which both started from juxtaposed clauses with independent speech acts, but ended up as complementizers after verbs denoting manipulative speech acts (`say'/`order'/`request' that p). Compare also, in general, the evolution of subordinate \textit{da}-clauses in South Slavic. For a comprehensive treatment cf.\ \citeauthor{Wiemer2017} (\citeyear[297--300, 313--315]{Wiemer2017}, \citeyear{Wiemer2019}).




\subsection{Preventives: The interplay between aspect and negation}
\label{subsec:preventives}


\textsc{NegDir} strategies are of an entirely different provenance. They strongly depend on the grammatical system, since their interpretation is conditioned by the key opposition of perfective vs.\ imperfective verb stems. The preventive function of negated perfective imperatives follows from the continuous entrenchment and inner-Slavic differentiation of the aspect system, which is based on stem-derivation and therefore applies to any finite or non-finite form of the verb. As a consequence, the choice between members of a binary opposition (\textsc{pfv} vs.\ \textsc{ipfv} stem) is unavoidable (\citealt{Wiemer2008}, \citealt{Seržant2017}). A preventive function of negated imperatives could probably always be inferred under favourable discourse conditions, just as in English or German today.\footnote{Compare Engl. \textit{Don't go there!}, Germ. \textit{Geh da nicht hin!} Such utterances can be understood either as prohibitives (`because I don't want you to') or as preventives (`because otherwise undesirable consequences will follow'), possibly apart from other readings discussed in \sectref{sec:negdir}.} However, the innovation (\citealt[29--30]{Ivić1958} and passim) consists in an increase of reliability of the choice between perfective and imperfective verb in the negated imperative, whereby perfective verbs in this grammatical context came to be more and more restricted to preventive meaning (as opposed to imperfective verbs becoming associated with prohibitive meaning).



Rather casual observations on this issue (\citealt{Wiemer2008}: 396--397) are corroborated by systematic corpus studies of early stages of South and East Slavic. Thus, \citet{MacRobert2013} points out that, although the admonitory use of the negated perfective imperative was ``well attested in Old Church Slavonic'' (\citeyear[286]{MacRobert2013}), its distribution was not identical with the situation in modern Slavic languages. Notably, other uses of the negated perfective imperative were all characterized by a focus on boundaries (especially in appeals to God to not do what He could do, using His might, at the present moment). The focus on boundaries is a common denominator of using perfective verbs, also in the imperative (see \sectref{subsec:negimp}), but obviously the variety of functions and discourse constellations used to be broader, before it shrank in North Slavic to leave the preventive function. Furthermore, \citeauthor{Mišina2020a} (\citeyear{Mišina2020b}, \citeyear{Mišina2020a}) carried out comprehensive corpus studies of early East Slavic texts and compared the frequency of negated imperatives in these texts with those in modern Russian: the proportion between imperfective vs.\ perfective stems for this grammatical form has remained more or less stable (3:1), however in early East Slavic the distribution between prohibitive and preventive function was much less clear-cut than in modern Russian, partly due to the elimination of perfective verbs from prohibitive contexts (\citeyear [172--176]{Mišina2020b}). Therefore, the crucial difference from South Slavic is that perfective verbs have not been driven out altogether from negated imperatives (see below), but have found a niche in the preventive function. This process can therefore be captured as pragmatic strengthening, not of particular markers, but of a grammatical distinction as such, and this process consisted in an increasing division of labour in the choice of perfective vs.\ imperfective stems (\sectref{subsec:negimp}). As far as we know by now, this innovation proves most consistent in the northeastern part of Slavic, but also the West Slavic languages seem to be rather consistent in this respect (see \sectref{subsec:areal}).


South Slavic languages, largely with the exception of Slovene, have gone another way: perfective stems have altogether been abandoned from negated imperatives (apart from petrified phrases). Already in many Old Church Slavonic monuments\footnote{Old Church Slavonic was based on east South Slavic dialects.} we can discern a preponderance of imperfective stems for negated imperatives regardless of their function (\citealt[71]{Kurzova-Ribarova1972}, \citealt[70--73]{Večerka1996}). The recent situation in South Slavic looks like a consequent continuation of this trend, accompanied by the spread of periphrastic constructions building on \textit{da}-clauses and ``analytical imperatives'' (with auxiliaries, Serbian/Macedonian \textit{nemoj}(\textit{te}), Bulgarian \textit{nedej}(\textit{te})). Among these newer constructions only the Balkan Slavic independent \textit{da}-clauses (with present tense forms) under negation seem to show a more or less reliable functional distribution between perfective (${\rightarrow}$ preventive) and imperfective (${\rightarrow}$ prohibitive) aspect, while analytical constructions with auxiliaries do not discriminate very consistently between these two functions (\sectref{subsec:analyticnegimp}). The \textit{da}-clauses, in turn, do not clearly discriminate between preventive and negative-optative or negative-purpose function (see \sectref{sec:negopt}).

The ousting of negated perfective imperatives by periphrastic constructions probably started from the (south)eastern end (i.e.\ Bulgarian); cf.\ \citet{Ivić1958}. Slovene, which is situated at the opposite end of the South Slavic continuum, has been least affected by this process; perfective verbs are used in negated imperatives, particularly with preventive meaning, but the distribution in opposition to imperfective verbs (${\rightarrow}$ prohibitive meaning) is arguably less clear-cut than in Russian. In this sense, it is more conservative.



\section{Summary and outlook}
\label{sec:summary}

The \textsc{NegDir} strategies are more homogeneous than the \textsc{NegOpt} strategies, even if among the latter we disregard patterns with an affinity to (or provenance from) negative purpose clauses. Functionally, the homogeneity of the \textsc{NegDir} strategies derives from the interaction of aspect with negation: either negated perfective imperatives (known as preventives) in North Slavic and Slovene or their \textit{da}-headed perfective present-tense equivalents in South Slavic (often considered ``analytical subjunctives"). However, the Balkan Slavic combination \textit{da ne} (\textsc{irr+neg}) reveals a high degree of morphological and semantic opacity exactly in apprehensional (including preventive) usage, while the preventive imperatives in North Slavic (and Slovene) are entirely compositional. Concomitantly, the aspect-based division of negated directive speech acts has grown very reliable in North Slavic, whereas South Slavic (with the exception of its northwestern periphery) has ousted negated perfective imperatives, and more recently developed periphrastic imperatives (with auxiliaries, Serbian/Macedonian \textit{nemoj(te)}, Bulgarian \textit{nedej(te)}) do not reliably distinguish between prohibitive and preventive function, even if they allow for perfective verbs.

Moreover, preventives in North Slavic are always distinguishable from the \textsc{Neg}\textsc{Opt} strategy, which is formally marked by clause-initial connectives. In South Slavic, particularly in Balkan Slavic, the two strategies are considerably less clearly distinguishable, both in functional and in structural terms: \textit{da ne}-clauses cover both the \textsc{NegOpt} and the \textsc{NegDir} function (although the former seems to be more salient), apart from other usage types. In addition, \textit{da ne} marks negative purpose clauses, but as a unit that is transparently composed of \textit{da} as a general irrealis morpheme (able to mark also purpose) and the negation.

\textsc{NegOpt} markers are based on irrealis markers with a very clear areal division: South Slavic has \textit{da}, North Slavic has \textit{by} (with Bulgarian \textit{da ne bi} showing morphological coalescence of both etyma). In \textsc{NegOpt} connectives they have however lost their morpheme status; in South Slavic, but not in North Slavic, we can furthermore observe a loss of TAM-constraints. The degree to which \textsc{NegOpt} connectives are etymologically and functionally related to optative or purpose clauses varies. Again, for the North Slavic languages these two sources of apprehensional clauses can be discerned more easily, together with areal tendencies: in Russian, negative purpose clauses (\textit{čtoby $\ldots$ ne}) have largely become dissociated from apprehensional semantics, while in West Slavic we observe an overlap between negative purpose and apprehensional complement clauses. Moreover, modern Russian demonstrates a tendency toward a syntactic division of labor between \textit{kak by $\ldots$  ne}, which is more frequent in complement than in independent clauses, and \textsc{NegOpt} markers based on particles, which are only used in independent clauses. Among other things, West Slavic is understudied particularly in this respect. In turn, the western part of South Slavic shows a tendency for \textsc{NegOpt} connectives to drop the requirement for clausal negation, together with a loss of TAM-constraints. Finally, the western part of Slavic (Slovene, Czech, Slovak) disprefers or altogether avoids infinitival complements, so that in apprehensional complement clauses same-subject and different-subject clause combining cannot be distinguished. By contrast, Russian consistently employs the finite-infinitive distinction, while Polish in this respect occupies an intermediary position.


Syntactically independent employment of apprehensional clauses can be considered a result of insubordination only if such clauses diachronically derive from negative purpose clauses. Only for Russian \textit{čtoby $\ldots$  ne} (not for \textit{kak by $\ldots$  ne}) is there sufficient evidence in favour of this assumption, while neither for West Slavic (Polish \textit{żeby $\ldots$  nie}, Polish/Czech/Slovak \textit{aby $\ldots$ ne}) nor for South Slavic \textit{da ne}-clauses any clear evidence exists, and there remains a plausible alternative hypothesis that negative purpose and apprehensional clauses developed in parallel. For all other strategies that are associated closely with optative semantics, insubordination is not an issue at all.

In conclusion, I want to emphasize that here I have only tried to give a comprehensive account of two types of strategies which, in my view, represent the most salient and widespread techniques across Slavic languages. However, in order to acquire a better understanding of how such strategies emerge and may become established in speaker communities, two directions of research urgently need to be pursued. First, we need more intensive research in the diachronic background of apprehensional constructions and their relation to other clause types. Second, on the contemporary level research should concentrate on the discovery of discourse-conditioned patterns which favour the emergence of apprehensional strategies.

\section*{Abbreviations}
\begin{tabularx}{.45\textwidth}{lQ}
$\cong$ & precedes equivalence of meaning \\
1/2/3 & first/second/third person \\
\textsc{acc} & accusative \\
\textsc{adv} & adverb \\
\textsc{aor} & aorist \\
\textsc{aux} & auxiliary \\
\textsc{comp} & complementizer \\
\textsc{con} & connective (unspecific) \\
CTP & complement-taking predicate \\
\textsc{cvb} & converb \\
\textsc{dat} & dative \\
\textsc{def} & definite article \\
\textsc{do} & definite direct object (Bulgarian masculine nouns) \\
DS & different subject \\
\textsc{excl} & exclamation\end{tabularx}%
\begin{tabularx}{.45\textwidth}{lQ}
\textsc{f} & feminine \\
\textsc{fut} & future \\
\textsc{gen} & genitive \\
\textsc{imp} & imperative \\
\textsc{impf} & imperfect \\
\textsc{ins} & instrumental \\
\textsc{ipfv} & imperfective \\
\textsc{irr} & irrealis \\
\textsc{lf} & \textit{l}-form (< participle in compound tenses) \\
\textsc{loc} & locative \\
\textsc{m} & masculine \\
\textsc{n} & neuter\\
\textsc{neg} & negation \\
\textsc{nom} & nominative \\
\textsc{pfv} & perfective \\
\textsc{pl} & plural \\
\textsc{pn} & proper name \\
\textsc{poss} & possessive \end{tabularx}%

\begin{tabularx}{.45\textwidth}{lQ}
\textsc{pp} & past (= anteriority) participle \\
\textsc{prs} & present \\
\textsc{pst} & past \\
\textsc{ptc} & particle\\
\textsc{q} & interrogative (= indefinite) pronoun  \end{tabularx}%
\begin{tabularx}{.45\textwidth}{lQ}
\textsc{refl} & reflexive marker \\ 
\textsc{sbjv} & subjunctive \\
\textsc{sg} & singular \\
SS & same subject \\
\textsc{voc} & vocative \\
\\\\
\end{tabularx}

\section*{Acknowledgements}
I am indebted to Eleni Bužarovska, Elena Ivanova and, in particular, to Nina Dobrushina and Eva Schultze-Berndt for useful comments based on an attentive reading of a previous version of this paper as well as to two anonymous reviewers for their helpful remarks. I also thank Jasmina Grković-Major, Ekaterina Mišina, Liljana Mitkovska, and Mladen Uhlik for their consultations, and I appreciate the exchange of information with Adriaan Barentsen, Tilman Berger, Markus Giger, Marc Greenberg, Leonid Iomdin, Maksim Makartsev, Alexander Rostovtsev-Popiel, Sergej Say, Bohumil Vykypěl, Natalia Zaika, and Nadežda Zorixina-Nilsson. Of course, the usual disclaimers apply.


\printbibliography[heading=subbibliography,notkeyword=this]
\noindent \Large\textbf{Corpora} \\

\normalsize

\noindent PNC:	Polish National Corpus, http://nkjp.pl/

\noindent RNC:	Russian National Corpus, http://www.ruscorpora.ru/new/

\end{document}
